\newcommand{\customDir}{../}

\RequirePackage{ifthen,xifthen}

% Input inkl. Umlaute, Silbentrennung
\RequirePackage[T1]{fontenc}
\RequirePackage[utf8]{inputenc}

% Arbeitsordner (in Abhängigkeit vom Master) Standard: .LateX_master Ordner liegt im Eltern-Ordner
\providecommand{\customDir}{../}
\newcommand{\setCustomDir}[1]{\renewcommand{\customDir}{#1}}
%%% alle Optionen:
% Doppelseitig (mit Rand an der Innenseite)
\newboolean{twosided}
\setboolean{twosided}{false}
% Eigene Dokument-Klasse (alle KOMA möglich; cheatsheet für Spicker [3 Spalten pro Seite, alles kleiner])
\newcommand{\customDocumentClass}{scrreprt}
\newcommand{\setCustomDocumentClass}[1]{\renewcommand{\customDocumentClass}{#1}}
% Unterscheidung verschiedener Designs: htw, fjs
\newcommand{\customDesign}{htw}
\newcommand{\setCustomDesign}[1]{\renewcommand{\customDesign}{#1}}
% Dokumenten Metadaten
\newcommand{\customTitle}{}
\newcommand{\setCustomTitle}[1]{\renewcommand{\customTitle}{#1}}
\newcommand{\customSubtitle}{}
\newcommand{\setCustomSubtitle}[1]{\renewcommand{\customSubtitle}{#1}}
\newcommand{\customAuthor}{}
\newcommand{\setCustomAuthor}[1]{\renewcommand{\customAuthor}{#1}}
%	Notiz auf der Titelseite (A: vor Autor, B: nach Autor)
\newcommand{\customNoteA}{}
\newcommand{\setCustomNoteA}[1]{\renewcommand{\customNoteA}{#1}}
\newcommand{\customNoteB}{}
\newcommand{\setCustomNoteB}[1]{\renewcommand{\customNoteB}{#1}}
% Format der Signatur in Fußzeile:
\newcommand{\customSignature}{\ifthenelse{\equal{\customAuthor}{}} {} {\footnotesize{\textcolor{darkgray}{Mitschrift von\\ \customAuthor}}}}
\newcommand{\setCustomSignature}[1]{\renewcommand{\customSignature}{#1}}
% Format des Autors auf dem Titelblatt:
\newcommand{\customTitleAuthor}[1]{\textcolor{darkgray}{Mitschrift von #1}}
\newcommand{\setCustomTitleAuthor}[1]{\renewcommand{\customTitleAuthor}{#1}}
% Standard Sprache
\newcommand{\customDefaultLanguage}[1]{}
\newcommand{\setCustomDefaultLanguage}[1]{\renewcommand{\customDefaultLanguage}{#1}}
% Folien-Pfad (inkl. Dateiname ohne Endung und ggf. ohne Nummerierung)
\newcommand{\customSlidePath}{}
\newcommand{\setCustomSlidePath}[1]{\renewcommand{\customSlidePath}{#1}}
\setCustomTitle{Angewandte Sensorik}
\setCustomSubtitle{Vorlesungsskript}
\setCustomAuthor{Falk-Jonatan Strube}
\setCustomNoteA{oder: „Vodka Gorbunoff: Der Sensorik reinster Seele“}
\setCustomNoteB{Vorlesung von Prof. Dr.-Ing. habil. Gorbunoff}
\RequirePackage{ifthen,xifthen}
%%% alle Optionen:
% Doppelseitig (mit Rand an der Innenseite)
\newboolean{twosided}
\setboolean{twosided}{false}
% Eigene Dokument-Klasse (alle KOMA möglich; cheatsheet für Spicker [3 Spalten pro Seite, alles kleiner])
\newcommand{\customDocumentClass}{scrreprt}
% Unterscheidung verschiedener Designs: htw, fjs
\newcommand{\customDesign}{htw}

%-- Für Spicker einiges anders:
\ifthenelse{\equal{\customDocumentClass}{cheatsheet}}{
	\documentclass[a4paper,10pt,landscape]{scrartcl}
	\usepackage{geometry}
	\geometry{a4paper, top=2mm, bottom=2mm, headsep=0mm, footskip=0mm, left=2mm, right=2mm}
	
	% Für Spicker \spsection für Section, zur Strukturierung \HRule oder \HDRule Linie einsetzen
	\usepackage{multicol}
	\newcommand{\spsection}[1]{\textbf{#1}}
	\setlist{nolistsep}
	\usepackage{parskip}
	\setlength{\parskip}{0em}
	\newcommand{\HRule}[1][\medskipamount]{\par
	  \vspace*{\dimexpr-\parskip-\baselineskip+#1}
	  \noindent\rule[0.2ex]{\linewidth}{0.2mm}\par
	  \vspace*{\dimexpr-\parskip-.5\baselineskip+#1}}
	\usepackage{dashrule}
	\newcommand{\HDRule}[1][\medskipamount]{\par
	  \vspace*{\dimexpr-\parskip-\baselineskip+#1}
	  \noindent\hdashrule[0.2ex]{\linewidth}{0.2mm}{1mm} \par
	  \vspace*{\dimexpr-\parskip-.5\baselineskip+#1}}
}{	%-- Ende Spicker-Unterscheidung-if
	%-- Unterscheidung Doppelseitig
	\ifthenelse{\boolean{twosided}}{
		\documentclass[a4paper,11pt, footheight=26pt,twoside]{scrreprt}
		\usepackage[head=23pt]{geometry}	% head=23pt umgeht Fehlerwarnung, dafür größeres "top" in geometry
		\geometry{top=30mm, bottom=22mm, headsep=10mm, footskip=12mm, left=20mm, right=20mm}
	}{
		\documentclass[a4paper,11pt, footheight=26pt]{scrreprt}
		\usepackage[head=23pt]{geometry}	% head=23pt umgeht Fehlerwarnung, dafür größeres "top" in geometry
		\geometry{top=30mm, bottom=22mm, headsep=10mm, footskip=12mm, inner=27mm, outer=13mm}
	}
	
	%-- Nummerierung bis Subsubsection für Report
	\ifthenelse{\equal{\customDocumentClass}{report} \OR \equal{\customDocumentClass}{scrreprt}}{
		\setcounter{secnumdepth}{3}	% zählt auch subsubsection
		\setcounter{tocdepth}{3}	% Inhaltsverzeichnis bis in subsubsection
	}{}
	
	\usepackage[perpage]{footmisc}	% Fußnotenzählung Seitenweit, nicht Dokumentenweit
	\usepackage{scrlayer-scrpage}	% Kopf-/Fußzeile

} %-- Ende Spicker-Unterscheidung-else

% Input inkl. Umlaute, Silbentrennung
\usepackage[T1]{fontenc}
\usepackage[utf8]{inputenc}
%%% if standard-sprache
\usepackage[english,ngerman]{babel}
\usepackage{csquotes}	% Anführungszeichen
\usepackage{eurosym}

% Style-Aufhübschung
\usepackage{soul, color}	% Kapitälchen, Unterstrichen, Durchgestrichen usw. im Text
%\usepackage{titleref}
\renewcommand*{\thefootnote}{\fnsymbol{footnote}}	% Fußnoten-Symbole anstatt Zahlen
\renewcommand*{\titlepagestyle}{empty} % Keine Seitennummer auf Titelseite

% Mathe usw.
\usepackage{amssymb}
\usepackage{amsthm}
\usepackage[fleqn,intlimits]{amsmath}	% fleqn: align-Umgebung rechtsbündig; intlimits: Integralgrenzen immer ober-/unterhalb
%\usepackage{mathtools} % u.a. schönere underbraces
\usepackage{xcolor}
\usepackage{esint}	% Schönere Integrale, \oiint vorhanden
\everymath=\expandafter{\the\everymath\displaystyle}	% Mathe Inhalte werden weniger verkleinert
\usepackage{wasysym}	% mehr Symbole, bspw \lightning
% Auch arcus-Hyperbolicus-Funktionen
\DeclareMathOperator{\arccot}{arccot}
\DeclareMathOperator{\arccosh}{arccosh}
\DeclareMathOperator{\arcsinh}{arcsinh}
\DeclareMathOperator{\arctanh}{arctanh}
\DeclareMathOperator{\arccoth}{arccoth} 
% Mathe in Anführungszeichen:
\newsavebox{\mathbox}\newsavebox{\mathquote}
\makeatletter
\newcommand{\mq}[1]{% \mathquotes{<stuff>}
  \savebox{\mathquote}{\text{"}}% Save quotes
  \savebox{\mathbox}{$\displaystyle #1$}% Save <stuff>
  \raisebox{\dimexpr\ht\mathbox-\ht\mathquote\relax}{"}#1\raisebox{\dimexpr\ht\mathbox-\ht\mathquote\relax}{''}
}
\makeatother
%\renewcommand{\int}{\int\limits}

% tikz usw.
\usepackage{tikz}
\usepackage{pgfplots}
\pgfplotsset{compat=1.11}	% Umgeht Fehlermeldung
\usetikzlibrary{graphs}
%\usetikzlibrary{through}	% ???
\usetikzlibrary{arrows}
\usetikzlibrary{arrows.meta}	% Pfeile verändern / vergrößern: \draw[-{>[scale=1.5]}] (-3,5) -> (-3,3);
\usetikzlibrary{automata,positioning} % Zeilenumbruch im Node node[align=center] {Text\\nächste Zeile} automata für Graphen
\usetikzlibrary{matrix}
\usetikzlibrary{patterns}	% Schraffierte Füllung
\tikzstyle{reverseclip}=[insert path={	% Inverser Clip \clip
	(current page.north east) --
	(current page.south east) --
	(current page.south west) --
	(current page.north west) --
	(current page.north east)}
% Nutzen: 
%\begin{tikzpicture}[remember picture]
%\begin{scope}
%\begin{pgfinterruptboundingbox}
%\draw [clip] DIE FLÄCHE, IN DER OBJEKT NICHT ERSCHEINEN SOLL [reverseclip];
%\end{pgfinterruptboundingbox}
%\draw DAS OBJEKT;
%\end{scope}
%\end{tikzpicture}
]	% Achtung: dafür muss doppelt kompliert werden!
\usepackage{graphpap}	% Grid für Graphen
\tikzset{every state/.style={inner sep=2pt, minimum size=2em}}

% Tabular
\usepackage{longtable}	% Große Tabellen über mehrere Seiten
\usepackage{multirow}	% Multirow/-column: \multirow{2[Anzahl der Zeilen]}{*[Format]}{Test[Inhalt]} oder \multicolumn{7[Anzahl der Reihen]}{|c|[Format]}{Test2[Inhalt]}
\renewcommand{\arraystretch}{1.3} % Tabellenlinien nicht zu dicht
\usepackage{colortbl}
\arrayrulecolor{gray}	% heller Tabellenlinien
\usepackage{array}	% für folgende 3 Zeilen (für Spalten fester breite mit entsprechender Ausrichtung):
\newcolumntype{L}[1]{>{\raggedright\let\newline\\\arraybackslash\hspace{0pt}}m{\dimexpr#1\columnwidth-2\tabcolsep-1.5\arrayrulewidth}}
\newcolumntype{C}[1]{>{\centering\let\newline\\\arraybackslash\hspace{0pt}}m{\dimexpr#1\columnwidth-2\tabcolsep-1.5\arrayrulewidth}}
\newcolumntype{R}[1]{>{\raggedleft\let\newline\\\arraybackslash\hspace{0pt}}m{\dimexpr#1\columnwidth-2\tabcolsep-1.5\arrayrulewidth}}
\usepackage{caption}	% Um auch unbeschriftete Captions mit \caption* zu machen

% Nützliches
\usepackage{verbatim}	% u.a. zum auskommentieren via \begin{comment} \end{comment}
\usepackage{tabto}	% Tabs: /tab zum nächsten Tab oder /tabto{.5 \CurrentLineWidth} zur Stelle in der Linie
\NumTabs{6}	% Anzahl von Tabs pro Zeile zum springen
\usepackage{listings} % Source-Code mit Tabs
\usepackage{lstautogobble} 
\usepackage{enumitem}	% Anpassung der enumerates
\setlist[enumerate,1]{label=\arabic*.)}	% global andere Enum-Items
\newenvironment{anumerate}{\begin{enumerate}[label=\alph*.)]}{\end{enumerate}} % Alphabetische Aufzählung
\renewcommand{\labelitemiii}{$\scriptscriptstyle ^\blacklozenge$} % global andere 3. Item-Aufzählungszeichen
\usepackage{letltxmacro} % neue Definiton von Grundbefehlen
% Nutzen:
%\LetLtxMacro{\oldemph}{\emph}
%\renewcommand{\emph}[1]{\oldemph{#1}}

% Einrichtung von lst
\lstset{
basicstyle=\ttfamily, 
mathescape=true, 
%escapeinside=^^, 
autogobble, 
tabsize=2,
basicstyle=\footnotesize\sffamily\color{black},
frame=single,
rulecolor=\color{lightgray},
numbers=left,
numbersep=5pt,
numberstyle=\tiny\color{gray},
commentstyle=\color{gray},
keywordstyle=\color{green},
stringstyle=\color{orange},
morecomment=[l][\color{magenta}]{\#}
showspaces=false,
showstringspaces=false,
breaklines=true,
literate=%
    {Ö}{{\"O}}1
    {Ä}{{\"A}}1
    {Ü}{{\"U}}1
    {ß}{{\ss}}1
    {ü}{{\"u}}1
    {ä}{{\"a}}1
    {ö}{{\"o}}1
    {~}{{\textasciitilde}}1
}
\usepackage{scrhack} % Fehler umgehen
\def\ContinueLineNumber{\lstset{firstnumber=last}} % vor lstlisting. Zum wechsel zum nicht-kontinuierlichen muss wieder \StartLineAt1 eingegeben werden
\def\StartLineAt#1{\lstset{firstnumber=#1}} % vor lstlisting \StartLineAt30 eingeben, um bei Zeile 30 zu starten
\let\numberLineAt\StartLineAt

% BibTeX
\usepackage[backend=bibtex, bibencoding=ascii,
%style=authortitle, citestyle=authortitle-ibid,
%doi=false,
%isbn=false,
%url=false
]{biblatex}	% BibTeX
\usepackage{makeidx}
%\makeglossary
%\makeindex

% Grafiken
\usepackage{graphicx}
\usepackage{epstopdf}	% eps-Vektorgrafiken einfügen

% pdf-Setup
\usepackage{pdfpages}
\usepackage[bookmarks,%
bookmarksopen=false,% Klappt die Bookmarks in Acrobat aus
colorlinks=true,%
linkcolor=black,%
citecolor=red,%
urlcolor=green,%
]{hyperref}

%-- Unterscheidung des Stils
\ifthenelse{\equal{\customDesign}{htw}}{
	% HTW Corporate Design: Arial (Helvetica)
	\usepackage{helvet}
	\renewcommand{\familydefault}{\sfdefault}
}{}

% Titel, Autor usw. werden vor dem Anfang des Dokuments in einem Rutsch definiert…
\newcommand{\DTitel}[1]{\newcommand{\Dokumententitel}{#1}}
\newcommand{\DUntertitel}[1]{\newcommand{\Dokumentenuntertitel}{#1}}
\newcommand{\DAutor}[1]{\newcommand{\Dokumentenautor}{#1}}
\newcommand{\DNotiz}[1]{\newcommand{\Dokumentennotiz}{#1}}
\newcommand{\DSign}[1]{\newcommand{\Dokumentensignatur}{#1}}
\DSign{\footnotesize{\textcolor{darkgray}{Mitschrift von\\ \Dokumentenautor}}}
\newcommand{\Autorformat}[1]{\textcolor{darkgray}{Mitschrift von #1}}
\newcommand{\workingdir}{../}	% Arbeitsordner (in Abhängigkeit vom Master) Standard: LateX_master Ordner liegt im Eltern-Ordner
% … Deswegen folgendes erst Nach Dokumentenbeginn ausführen:
\AtBeginDocument{
	\hypersetup{
		pdfauthor={\Dokumentenautor},
		pdftitle={HTW Dresden | \Dokumententitel - \Dokumentenuntertitel},
	}
	\automark[section]{section}
	\automark*[subsection]{subsection}
	\pagestyle{scrheadings}
	\ihead{\includegraphics[height=1.7em]{\workingdir .LaTeX_master/HTW-Logo.eps}}
	%\ohead{\truncate{5cm}{\Dokumententitel}}
	\ohead{\Dokumententitel}
	\cfoot{\pagemark}
	\ofoot{\Dokumentensignatur}
	% Titelseite
	\title{\includegraphics[width=0.35\textwidth]{\workingdir .LaTeX_master/HTW-Logo.eps}\\\vspace{0.5em}
	\Huge\textbf{\Dokumententitel} \\\vspace*{0,5cm}
	\Large \Dokumentenuntertitel \\\vspace*{4cm}}
	\author{\Autorformat{\Dokumentenautor} \vspace*{1cm}\\\Dokumentennotiz}
}

% Unterabschnitte
%\newtheorem{example}{Beispiel}%[section]
%\newtheorem{definition}{Definition}[section]
%\newtheorem{discussion}{Diskussion}[section]
%\newtheorem{remark}{Bemerkung}[section]
%\newtheorem{proof}{Beweis}[section]
%\newtheorem{notation}{Schreibweise}[section]
\RequirePackage{xcolor}
%% EINFACHE BEFEHLE

% Abkürzungen Mathe
\newcommand{\EE}{\mathbb{E}}
\newcommand{\QQ}{\mathbb{Q}}
\newcommand{\RR}{\mathbb{R}}
\newcommand{\CC}{\mathbb{C}}
\newcommand{\NN}{\mathbb{N}}
\newcommand{\ZZ}{\mathbb{Z}}
\newcommand{\PP}{\mathbb{P}}
\renewcommand{\SS}{\mathbb{S}}
\newcommand{\cA}{\mathcal{A}}
\newcommand{\cB}{\mathcal{B}}
\newcommand{\cC}{\mathcal{C}}
\newcommand{\cD}{\mathcal{D}}
\newcommand{\cE}{\mathcal{E}}
\newcommand{\cF}{\mathcal{F}}
\newcommand{\cG}{\mathcal{G}}
\newcommand{\cH}{\mathcal{H}}
\newcommand{\cI}{\mathcal{I}}
\newcommand{\cJ}{\mathcal{J}}
\newcommand{\cM}{\mathcal{M}}
\newcommand{\cN}{\mathcal{N}}
\newcommand{\cP}{\mathcal{P}}
\newcommand{\cR}{\mathcal{R}}
\newcommand{\cS}{\mathcal{S}}
\newcommand{\cZ}{\mathcal{Z}}
\newcommand{\cL}{\mathcal{L}}
\newcommand{\cT}{\mathcal{T}}
\newcommand{\cU}{\mathcal{U}}
\newcommand{\cV}{\mathcal{V}}
\renewcommand{\phi}{\varphi}
\renewcommand{\epsilon}{\varepsilon}

% Farbdefinitionen
\definecolor{red}{RGB}{180,0,0}
\definecolor{green}{RGB}{75,160,0}
\definecolor{blue}{RGB}{0,75,200}
\definecolor{orange}{RGB}{255,128,0}
\definecolor{yellow}{RGB}{255,245,0}
\definecolor{purple}{RGB}{75,0,160}
\definecolor{cyan}{RGB}{0,160,160}
\definecolor{brown}{RGB}{120,60,10}

\definecolor{itteny}{RGB}{244,229,0}
\definecolor{ittenyo}{RGB}{253,198,11}
\definecolor{itteno}{RGB}{241,142,28}
\definecolor{ittenor}{RGB}{234,98,31}
\definecolor{ittenr}{RGB}{227,35,34}
\definecolor{ittenrp}{RGB}{196,3,125}
\definecolor{ittenp}{RGB}{109,57,139}
\definecolor{ittenpb}{RGB}{68,78,153}
\definecolor{ittenb}{RGB}{42,113,176}
\definecolor{ittenbg}{RGB}{6,150,187}
\definecolor{itteng}{RGB}{0,142,91}
\definecolor{ittengy}{RGB}{140,187,38}

% Textfarbe ändern
\newcommand{\tred}[1]{\textcolor{red}{#1}}
\newcommand{\tgreen}[1]{\textcolor{green}{#1}}
\newcommand{\tblue}[1]{\textcolor{blue}{#1}}
\newcommand{\torange}[1]{\textcolor{orange}{#1}}
\newcommand{\tyellow}[1]{\textcolor{yellow}{#1}}
\newcommand{\tpurple}[1]{\textcolor{purple}{#1}}
\newcommand{\tcyan}[1]{\textcolor{cyan}{#1}}
\newcommand{\tbrown}[1]{\textcolor{brown}{#1}}

% Umstellen der Tabellen Definition
\newcommand{\mpb}[1][.3]{\begin{minipage}{#1\textwidth}\vspace*{3pt}}
\newcommand{\mpe}{\vspace*{3pt}\end{minipage}}

\newcommand{\resultul}[1]{\underline{\underline{#1}}}
\newcommand{\parskp}{$ $\\}	% new line after paragraph
\newcommand{\corr}{\;\widehat{=}\;}
\newcommand{\mdeg}{^{\circ}}

\newcommand{\nok}[2]{\begin{pmatrix}#1\\#2\end{pmatrix}}	% n über k BESSER: \binom{n}{k}
\newcommand{\mtr}[1]{\begin{pmatrix}#1\end{pmatrix}}	% Matrix
\newcommand{\dtr}[1]{\begin{vmatrix}#1\end{vmatrix}}	% Determinante (Betragsmatrix)
\renewcommand{\vec}[1]{\underline{#1}}	% Vektorschreibweise
\newcommand{\imptnt}[1]{\colorbox{red!30}{#1}}	% Wichtiges
\newcommand{\intd}[1]{\,\mathrm{d}#1}
\newcommand{\diffd}[1]{\mathrm{d}#1}

%\documentclass[a4paper,18pt]{scrreprt}

\usepackage[ngerman]{babel}
\usepackage{fontspec} 
\usepackage{unicode-math}
\usepackage[a4paper]{geometry}
\geometry{verbose,tmargin=2cm,bmargin=2cm,lmargin=2cm,rmargin=2cm,headheight=1cm,headsep=1cm,footskip=1cm}
\usepackage{graphicx}

\providecommand{\customTitle}{Titel}
\providecommand{\customAuthor}{Autor}

\title{\customTitle}
\author{\customAuthor}

\usepackage{xcolor}
\definecolor{darkblue}{RGB}{0,0,102}
\color{darkblue}

\providecommand{\imptnt}[1]{#1}
\usepackage{amsmath}
\usepackage{enumitem}	% Anpassung der enumerates

\usepackage{geometry}
\geometry{top=10mm, bottom=15mm, left=15mm, right=15mm}

\setmainfont{FJS} 
\setsansfont{FJS} 
\setmonofont{FJS}

\setmathfont{FJS}
\setmathfont[range=\mathup]  {FJS}
\setmathfont[range=\mathbfup]{FJS}
\setmathfont[range=\mathbfit]{FJS}
\setmathfont[range=\mathit]  {FJS}

\usepackage{enumitem}	% Anpassung der enumerates

\usepackage{luacode,luatexbase} 

\begin{luacode}
local function replaceSomething ( line )
  return string.gsub(line, "\\Rightarrow" , "=>") 
end
luatexbase.add_to_callback( "process_input_buffer",  replaceSomething, "replace_something")
\end{luacode}

\usepackage{array}	% für folgende 3 Zeilen (für Spalten fester breite mit entsprechender Ausrichtung):
\newcolumntype{L}[1]{>{\raggedright\let\newline\\\arraybackslash\hspace{0pt}}m{\dimexpr#1\columnwidth-2\tabcolsep-1.5\arrayrulewidth}}
\newcolumntype{C}[1]{>{\centering\let\newline\\\arraybackslash\hspace{0pt}}m{\dimexpr#1\columnwidth-2\tabcolsep-1.5\arrayrulewidth}}
\newcolumntype{R}[1]{>{\raggedleft\let\newline\\\arraybackslash\hspace{0pt}}m{\dimexpr#1\columnwidth-2\tabcolsep-1.5\arrayrulewidth}}
\usepackage{caption}	% Um auch unbeschriftete Captions mit \caption* zu machen


\begin{document}

\maketitle
\newpage
\tableofcontents
\newpage

\chapter*{Vorbemerkung}
Prüfung: Keine Rechnung, kein Programmieren. Reine Wissensabfrage (Welche Sensoren für diesen Zweck, Vergleich, Vorteil/Nachteil, …)\\
Hilfsmittel: alles Handschriftliche zugelassen.

\setcounter{chapter}{-1}
\chapter{Einführung}

Elektronische Geräte (PC usw.) sind blind, brauchen Sensoren.

\section{Eingabegeräte}
Eingabegeräte: alle Geräte, über die einem Computer Informationen zugeführt werden können.
\begin{itemize}
\item \emph{Umfeldsensoren}: führen die Informationen über die Umgebung unmittelbar dem Computer zu (ohne menschliches einwirken): Temperatur, Druck, Schall, … $\Rightarrow$ Hauptthema der VL
\item \emph{User Interface Devices}: Informationseingabe durch Benutzer: Tastaturen, Touchscreens (kapazitiv, induktiv, …) usw.
\end{itemize}

\section{Ziel des Kurses}
\begin{itemize}
\item \emph{Sensorik} (Sensortechnik): Teilgebiet der Messtechnik, das sich mit Entwicklung und Einsatz von Sensoren befasst.
\item wichtigste physikalischen Prinzipien zur Erfassung nichtelektrischer Messgrößen
\item Funktionsweise ausgewählter Sensoren
\item Vor- und Nachteile jedes Typs
\item Anwendungshinweise für den zuverlässigen Sensorenbetrieb
\end{itemize}

Kapitel:
\begin{itemize}
\item Grundlagen der Messtechnik
\item Einführung in die Sensorik
\item Temperatursensoren
\item Drucksensoren
\item Durchflusssensoren
\item Binäre Positionssensoren
\end{itemize}

\chapter{Grundlagen der Messtechnik}

\section{Messgröße und Messwert}

\begin{itemize}
\item \imptnt{\emph{Messgröße}} ist ein messbares \imptnt{Merkmal} eines Objektes oder Prozess:
\begin{itemize}
\item Eigenschaft 
\item Vorgang
\item Zustand
\end{itemize}
$\Rightarrow$ Temperatur, Druck, Durchfluss, Position
\item \imptnt{\emph{Maßeinheit}}: eine durch Vereinbarung festgelegte \imptnt{Vergleichsgröße}
\item \emph{SI-System:} m, kg, s, A, K, mol, cd
\item \imptnt{\emph{Messwert}}: der gemessene Wert einer Messgröße 
$=ZW \cdot \mathrm{Einheit}$ 
%$=ZW * \mathrm{Einheit}$ 
(also Messgröße kombiniert mit Maßeinheit)
\item Messung: (Ablauf)
\begin{itemize}
\item Quelle (Träger eine Messgröße)
\item Erfassung der Größe \& Vergleich der Größe mit einer Einheit
\item Anzeige (des Messwerts)
\end{itemize}
\end{itemize}

\subsection{Normal}
\emph{Normal} (DIN 1390): Maßverkörperung, Messgerät, Referenzmaterial oder Messeinrichtung zum \linebreak 
Zweck eine Einheit festzulegen, zu verkörpern, zu bewahren oder zu reproduzieren.
\begin{itemize}
\item Internationale (nationale) Normale (Kopie der Ursprungs-Normale: bspw. Kopie Ur-Kilo)
\item Bezugsnormale (PTB) (für Betriebe usw.)
\item Gebrauchs-, Prüfnormale (für die konkrete Messung)
\item Lehren
\item …
\end{itemize}

\subsection{Messen}
\emph{Messen} ist der Vorgang, durch den ein Wert einer physikalischen Größe als Vielfaches eines Bezugswertes (Maßeinheit) ermittelt wird (Messen = Vergleichen).

\subsection{Messprinzipien}
\emph{Messprinzip}: eine charakteristische naturwissenschaftliche Erscheinung oder ein gesetzmäßiger Zusammenhang, der einer Messung zugrunde gelegt wird.

Physikalische Grundlage der Messung: über 100 physikalisch-chemisch-biologische Effekte.

\section{Messmethoden}
\emph{Messmethode}: „spezielle, vom Messprinzip unabhängige Art des Vorgehens bei der Messung.“ (DIN 1319-1)
\begin{itemize}
\item allgemeine Vorgehensweise bei der Durchführung von Messungen
\item nicht an eine physikalische Realisierung gebunden
\end{itemize}

\subsection{Methoden nach Art der Messgröße}
Bsp. Messung der Temperatur (die sich nicht ändert)
\begin{itemize}
\item \emph{statisch}: bestimmt wird eine zeitlich unveränderliche Messgröße nach einem Messprinzip, das nicht auf der zeitlichen Änderung anderer Größen beruht (Bsp.: Ablesen der Flüssigkeits-Säule im Thermometer)
\item \emph{dynamisch}: Messgröße ist entweder selbst zeitlich veränderlich oder ihr Wert ergibt sich aus zeitlichen Änderungen anderer Größen (Messung durch Einwirkung anderer Kräfte um die Reaktion heraus zu finden [und so zu messen]. Bsp. Halbwertszeit von Materialien lässt Rückschlüsse auf Temperatur zu)
\end{itemize}

\subsection{Methoden nach Art der Vergleichsgröße}
\begin{itemize}
\item \emph{direkte Messmethode}: der Wert der Messgröße wird durch quantitative Vergleich mit einer physikalischen gleichartigen Einheit ermittelt (Metermaß).
\item \emph{indirekte Messmethode}: der Messwert wird aus den Messdaten einer oder mehrerer anderer (direkt messbaren) Größen ermittelt, die mit der gesuchten Messgröße in einem definierten Zusammenhang steht (stehen) (Messung der Temperatur durch Messung der Ausdehnung der Flüssigkeit im Thermometer)
\item \emph{inkrementale Messmethode}: der Messwert wird von einem Bezugspunkt aus durch Addition/Subtraktion von kleinen Wertzuwachsen (Inkrementen) ermittelt (Geld zählen: Münzen (als Inkrement) in Säule stapeln, diese messen und damit den Wert bestimmen; vgl. auch Kugelmaus: gerasterte, inkrementale Messung der Veränderung der Kugelbewegung)
\end{itemize}

\subsection{Methoden nach Art der Vergleichsvorgehensweise}
\begin{itemize}
\item \emph{Ausschlagmethode}: Messgröße wird direkt oder über eine Zwischengröße in eine möglichst proportionale Ausgangsgröße umgewandelt (vglw. Küchenwaage mit Feder)\\
Messgröße $M$ $\to$ Übertragungsfunktion $f$ $\to$ Ausgangsgröße $A$\\
$\Rightarrow \boxed{M = f^{-1}(A)}$
\begin{itemize}[label=$+$]
\item keine Hilfsenergie benötigt
\item einfacher Aufbau
\item schnell
\end{itemize}
\begin{itemize}[label=-]
\item Energieentzug aus Messobjekt $\to$ Rückwirkung
\item Kennlinie (Übertragungsfunktion) des Messgerätes muss bekannt sein
\item starker Einfluss von Störgrößen
\end{itemize}
\item \emph{Kompensationsmethode}: Messgerät erzeugt eine bekannte variable Vergleichsgröße, um die Differenz zwischen beiden Größen gegen Null streben zu lassen.\\
ABB AS1\\
$M=S$\\
(Vergleich: alte Küchenwaage mit Schieberegler/Ausgleichsgewichten)
\begin{itemize}[label=$+$]
\item Variable Art der Kompensationsgröße
\item Reduzierung der Rückwirkung und Störeinflüsse
\item Nichtlinearität unkritisch
\item keine Kalibrierung notwendig
\item leichte Realisierung großer Messbereiche
\item hohe Genauigkeit möglich
\end{itemize}
\begin{itemize}[label=$-$]
\item größerer gerätetechnischer Aufwand
\item mehrstufig: in der Regel langsam
\item Hilfsenergie notwendig
\item Viele bekannte Vergleichsgrößen nötig
\end{itemize}
\end{itemize}
\paragraph{Beispiel Spannungsmessung}
\begin{itemize}
\item Ausschlagmethode\\
ABB AS2\\
Problem: braucht unendlich großen Widerstand, sonst verfälscht es die Messung
\item Kompensationsmethode\\
ABB AS3\\
Mit bekannter Spannung $U_{ref}$.\\
Strommessgerät: Ist Strom gleich 0? Dann sind beide Spannungen gleich.
\end{itemize}
weitere Methoden:
\begin{itemize}
\item Differenzierungsmethode
\item Substitutionsmethode
\end{itemize}

\subsection{Methoden nach Zeitablauf}
\begin{itemize}
\item \emph{kontinuierliche Messmethode}: die Größe wird ohne zeitliche Unterbrechung erfasst und auch angezeigt. (Vgl. Thermometer mit Flüssigkeit verändert sich stets/kontinuierlich)\\
ABB AS4
\item \emph{diskontinuierliche Messmethode}: die Messgröße wird nur zu bestimmten (diskreten) Zeitpunkten erfasst (abgetastet) und/oder angezeigt.\\
ABB AS5\\
Informationen werden begrenzt (aufs wesentliche), Zeit muss diskretisiert werden (in vernünftigen Maß: Nyquist-Shannon-Abtasttheorem $f_{Abtast}\geq f_{max}$)
\end{itemize}

\subsection{Methoden nach Wertebereich der Ausgangsgröße}
\begin{itemize}
\item \emph{analoge Messmethode}: die Messgröße wird durch eine eindeutige und stetige Anzeigegröße (Messwert) abgebildet.
\begin{itemize}[label=\textbullet]
\item kontinuierlicher Wertebereich
\end{itemize}
(Spannungsmesser mit Zeiger, Analoge Küchenwaage)
\item \emph{digitale Messmethode}: die Messgröße wird in Form einer in festgelegten Schritten quantisierten Anzeigegröße abgebildet.
\begin{itemize}[label=\textbullet]
\item diskreter Wertebereich
\end{itemize}
(digital Multimeter)
\end{itemize}

\section{Signale}
\begin{itemize}
\item \emph{Signal} in der Messtechnik: wenn man einer speziell ausgewählten \emph{zeitlich veränderlichen} physikalischen Größe (Signalträger, Informationsträger) eine Information zuordnet. \\
Hat 3 Elemente: Erzeuger, Träger, Empfänger (Kommunikation: Nachricht wird generiert (im Kopf) und Schall wird erzeugt (Sprechen) $\to$ Schall wird übertragen $\to$ Schall wird gehört (Ohr) und Schall-Information wird entschlüsselt (im Kopf))
\item \emph{Messsignal}: Signal mit einem (oder mehreren) zeitvariablen informationstragenden Parameter (Informationsparameter), der die Werte der Messgröße in eindeutiger und reproduzierbarer Weise abgebildet.
\item \emph{Rauschen}: zufällige zeitlich veränderliche Größe ohne nützliche Information.
\end{itemize}

\subsection{Informationsparameter}
Kontinuierlich:
\begin{itemize}
\item Signalträgerwert („Wert auf x-Achse“)
\item Amplitude
\item Frequenz
\item Phase
\end{itemize}
Diskret:
\begin{itemize}
\item Impulshöhe
\item Pulsbreite (Pulsbreite/-länge als Indikator für Wert)
\item Pulsfolge (Digitaltechnik: 001 hat andere Folge als 010)
\end{itemize}

\subsection{Signaltypen}
\begin{itemize}
\item \emph{Determiniertes Singal}\\
Der Signalwert ist zu jedem Zeitpunkt verfügbar
\begin{itemize}[label=$+$]
\item Information mit einmaliger Messung gewinnbar
\item[$-$] zu viel Information. Sie kann auch durch Störung unbrauchbar werden
\end{itemize}
\item Stochastisches Signal\\
regellos, zufällig schwankender Signalverlauf
\begin{itemize}[label=$+$]
\item Störungen machen sich nur stark reduziert bemerkbar, sie werden über die Messzeit integriert
\item[$-$] Information ist erst mit mehrmaligen Messungen zu gewinnen, das erfordert einen großen Zeitbedarf
\end{itemize}
\item \emph{Signalgemisch}\\
deterministisches Signal mit einem stochastischen Anteil (Rauschen)
\begin{itemize}
\item[$-$] Das Rauschen ist unerwünscht und muss unterdrückt oder ausgefiltert werden
\end{itemize}
\end{itemize}

\subsubsection{Nach Wertebereich}
\begin{itemize}
\item \emph{Wertkontinuierliches Signal}: kontinuierlicher Wertebereich
\begin{itemize}[label=$+$]
\item der Informationsparameter bildet adäquat die Messgröße ab
\item kann theoretisch beliebig viele Werte innerhalb seines Wertebereichs annehmen
\item[$-$] einfach zu stören
\end{itemize}
\item \emph{Wertdiskretes Signal}: diskontinuierlicher Wertebereich aus einer endlichen Anzahl von vordefinierten Werten
\begin{itemize}
\item[$+$] Störeinflüsse machen sich erst nach Überschreiten von Grenzwerten bemerkbar
\item[$-$] möglicher Informationsverlust (wenn Inkrement falsch gewählt wird)
\end{itemize}
\item \emph{binäres Signal}: der Wertebereich hat nur zwei Werte
\end{itemize}
\subsubsection{Nach zeitlichen Verlauf}
\begin{itemize}
\item \emph{kontinuierliches Signal}: IP kann zu jedem beliebigen Zeitpunkt seinen Wert ändern
\begin{itemize}[label=$+$]
\item jederzeit vorhanden: jederzeit ist der zeitliche Verlauf der Messwerten verfolgbar
\item[$-$] Störungen können jederzeit wirken
\item[$-$] Informationsmenge ist oft unnötig groß
\end{itemize}
\item \emph{diskontinuierliches Signal}: IP kann nur zu diskreten Zeitpunkten seinen wert ändern.
\begin{itemize}[label=$+$]
\item Störungen zwischen den Zeitpunkten der Parameterveränderungen können sich nicht auswirken
\item[$-$] Informationen stehen nur zu diskreten Zeitpunkten zur Verfügung (Informationsverlust)
\end{itemize}
\end{itemize}
\subsubsection{Kombinationen}
\begin{itemize}
\item Analogsignal
\item Digitalsignal (Zeit- und Wert-diskret)
\begin{itemize}[label=$+$]
\item hohe Störsicherheit
\item einfache Datenspeicherung
\item flexible Weiterverarbeitung
\item vielfältige Übertragungsmöglichkeiten
\end{itemize}
\item diskret-kontinuierlich ($\Delta A$: Amplituden-Quantisierungsintervall)
\item analog-diskontinuierlich ($\Delta t$: Zeit-Quantisierungsintervall)
\end{itemize}

\subsection{Normsignale}
Strom als Informationsträger ist störungssicher $\to$ Spannung ist leichter zu manipulieren als Stromstärke!
\begin{itemize}
\item \emph{Strom analog}:\\
$0\dots 20\; \mathrm{mA}$ (dead zero)\\
$4 \dots 20 \;\mathrm{mA}$ (live zero)
\item \emph{Spannung analog}:\\
$0\dots 10 \;\mathrm{mA}$ (dead zero)\\
$2 \dots 10 \;\mathrm{mA}$ (live zero)
\item \emph{Spannung digital}:\\
Spannungsbereiche an Transistoren für Definition von 1/0-Signalen (TTL, LVTTL, CMOS)
\end{itemize}
Unterscheidung dead/live zero um einen Kabelbruch von einem Null-Signal unterscheiden zu können. Bei normaler Temperatur Unterscheidung bspw. $0\;\mathrm{mA}$ dead zero und $4 \;\mathrm{mA}$ live zero.

\section{Messeinrichtung}
Gesamtheit aller zusammenhängender Funktionseinheiten (Glieder), die zum Zweck der Messung, Messdatenverarbeitung, Anzeige von einer oder mehreren Messgrößen benutzt werden [In Anlehnung an die Norm DIN 1319].

besteht aus:
\begin{itemize}
\item Messgeräte
\item Hilfsgeräte
\item Maßverkörperungen
\item chemische Reagenzien
\item …
\end{itemize}
Messanlage: dauerhaft installierte Messeinrichtung

\paragraph{Beispiel: Fieberthermometer}
\begin{itemize}
\item Eingangsgröße: Temperatur $\vartheta$
\item Eingang: Sensor Thermoelement PT100. Temperatur erzeugt Änderung des darin liegenden Widerstands 
\item Anpassung: Element erzeugt durch Spanngs-/Stromquelle vom Widerstand abhängige Signale
\item Verstärkung
\item AD-Wandler
\item $\mu$Controller (verarbeitet Signale; piepst, wenn Messung fertig)
\item Anzeige der Temperatur
\end{itemize}
\subsection{Messkette}
$\underbrace{\overset{\substack{\text{Mess-}\\\text{größe}}}{\longrightarrow} \text{ Aufnehmer (Sensor)}}_{\text{Prozess}}$ 
$\overset{\substack{\text{Mess-}\\\text{signal}}}{\longrightarrow}$ Messumformer 
$\overset{\substack{\text{Norm-}\\\text{signal}}}{\longrightarrow}$ Messgerät / Anzeige 
$\overset{\text{Ausgabe}}{\longrightarrow}$\\
offene Wirkungskette: Messkette

\section{Steuern}
Ein Vorgang in einem System, bei dem eine oder mehrere Größen als Eingangsgrößen (Führungsgrößen) andere Größen als Ausgangsgrößen aufgrund der dem System eigentümlichen Gesetzmäßigkeiten beeinflussen [DIN 19226].\\
Steuerung ist also ein Plan, wie eine veränderliche Größe beeinflusst werden soll.\\
Beispiel Steuerung: Wecker klingelt jeden Tag um 7 Uhr. In dem Fall steuert der Wecker die Person.\medskip

Stellglied (beeinflusst unmittelbar die Steuergröße) kann verstellt werden (Steuerglied+Steller erzeugen aus der Führungsgröße die \emph{Stellgröße}), um \emph{Steuergröße} (output) zu ändern. Der Sollwert (\emph{Führungsgröße}) für die Steuergröße wird vom \emph{Führungsglied} (als starrer Befehlsgröße) diktiert. Es gibt allerdings keine Rückmeldung der Steuergröße, \emph{Störgrößen} könnten also nicht bedacht werden.\\
$\underbrace{\text{Stellglied (Aktor)}}_{\substack{\text{Prozess}\\\text{(Steuergröße)}\\\text{über Steuerstrecke}}}$
$\overset{\substack{\text{Stell-}\\\text{größe}}}{\longleftarrow}$ Steller
$\underset{\substack{\text{(Norm-}\\\text{signal)}}}{\overset{\substack{\text{Stell-}\\\text{signal}}}{\longleftarrow}}$ Steuerglied 
$\overset{\substack{\text{Führungs-}\\\text{größe}}}{\longleftarrow}$ Führungsglied\\
offene Steuerkette\\
Nachteile:
\begin{itemize}
\item keine Rückmeldung
\item dynamische Eigenschaften des Prozesses müssen genau bekannt sein.
\end{itemize}

\section{Regeln}
Ein Vorgang, bei dem fortlaufend eine Größe (die zu regelnde Größe, die Regelgröße) erfasst, mit einer anderen Größe (der Führungsgröße) verglichen, und im Sinne einer Angleichung an die Führungsgröße beeinflusst wird [DIN 19226].
\begin{enumerate}
\item Messen: die Regelgröße wird direkt gemessen oder aus anderen Messgrößen berechnet.
\item Vergleichen: die Regelgröße wird mit der Führungsgröße verglichen und die Regelabweichung berechnet.
\item Regelglied + Steller erzeugt die Stellgröße aus der Regelabweichung
\item[\textbullet] mit der Stellgröße wird somit die Regelgröße wieder beeinflusst.
\item[\textbullet] eine Störgröße kann durch das Messen der Regelgröße wieder an den Steller zurück geführt werden.
\end{enumerate}
Aus dem Prozess (Regelgröße) führt somit eine offene Messkette als Ausgabe heraus und führt eine offene Steuerkette als Regelabweichung herein. So wird zwischen dem Ein- und Ausgang ein Vergleichsglied gesteckt, die aus der Rückführgröße (Ist-Wert) und der Führungsgröße (Soll-Wert, von Außen bestimmt) diese Regelabweichung berechnet. Im Prozess führt die Regelstrecke vom Eingang zum Ausgang.

So entsteht ein geschlossener \emph{Regelkreis}.

\section{Steuern oder Regeln?}
\paragraph{Steuerung}
\begin{itemize}
\item offene Wirkungskette
\item Prozesseigenschaften (statisch + dynamisch) müssen genau bekannt sein)
\item kann auf Störungen nicht reagieren
\item kein Soll-Ist-Vergleich
\item keine Messung notwending
\item Prozessstabilität wird nicht beeinflusst
\end{itemize}
ist sinnvoll, wenn:
\begin{itemize}
\item die Auswirkungen von Störgrößen vernachlässigbar klein sind
\item nur eine Störgröße auftritt, die nach Art und Verlauf bekannt ist
\item Störgrößenänderungen selten sind
\end{itemize}
\paragraph{Regelung}
\begin{itemize}
\item geschlossener Regelkreis
\item Prozesseigenschaften müssen nicht genau bekannt sein (Robustheit gegenüber Parameteränderungen)
\item kann Störungen ausregeln (Störkompensation)
\item Soll-Ist-Vergleich
\item Messung ist notwending
\item Der geschlossene Regelkreis kann instabil werden (bsp.: Mikrofon neben Lautsprecher $\to$ Pfeifen)
\end{itemize}
ist sinnvoll, wenn:
\begin{itemize}
\item veränderliche, unbekannte, nach Art und Größe verschiedene, nicht messbare Störgrößen auftreten können
\end{itemize}

\chapter{Sensorik}
\section{Sensor}
Der Sensor ist das primäre Element einer Messkette\\
Sensor = Aufnehmer, Messwertaufnehmer, Messwertgeber, Signalgeber, Fühler, Geber, Initiator, Detektor, Zelle, Wandler, Transducer\\
$\overset{\text{Messgröße}}{\longrightarrow} \underbrace{\underbrace{\underbrace{
\text{Aufnehmer}}_{\text{Sensorelement}}\overset{\text{Messsignal}}{\longrightarrow}
\text{Signalaufbereitung} \overset{\text{Normsignal}}{\longrightarrow}}_{\text{integrierter Sensor}}
\substack{\text{Signalverarbeitung}\\\text{und -auswertung}}\underset{\text{Steuersignal}}{\overset{\text{Ausgabe}}{\rightrightarrows}}}_{\text{intelligenter Sensor}}$
\section{Einteilung von Sensoren}
\begin{anumerate}
\item Funktionsweise
\begin{itemize}
\item selbstgenerierende Sensoren (keine externe Energiequelle)
\begin{itemize}
\item Spannungsausgang
\item Stromausgang
\item Ladungsausgang
\end{itemize}
\item modulierende Sensoren (externe Energiequelle nötig)
\begin{itemize}
\item Widerstand
\item Induktivität
\item Kapazität
\end{itemize}
\end{itemize}
\item Aktive und passive Sensoren
\begin{itemize}
\item Selbstgenerierende Sensoren\\
= aktive Sensoren (in dt. Literatur, in engl.: passive)
\item Modulierende Sensoren\\
= passive Sensoren (in dt. Literatur, in engl.: aktive)
\end{itemize}
Alternativ:
\begin{itemize}
\item passive: Aufnehmen vorhandener Signale\\
bspw. Kamera, Mikrofon, IR-Sensor
\item aktive: Stimulieren der Umwelt und Aufnehmen der Antwort\\
bspw. Ultraschallsensor, Laserscanner, Radar
\end{itemize}
\item Messprinzipien / Energieformen
\begin{itemize}
\item elektrisch
\item magnetisch
\item mechanisch
\item thermisch
\item Strahlung
\item chemisch/biologisch
\end{itemize}
\item Einzelne Messgrößen
\item Grad der Erfassung
\begin{itemize}
\item erfassende Sensoren (Binärsensoren)
\item messende Sensoren (Analog- und Digitalsensoren)
\end{itemize}
\item Kontaktierende und kontaktlose Sensoren
\item Absolute und relative Sensoren
\item Interne, externe Sensoren\\
typisch interner Sensor: Schmerz\\
typisch externer Sensor: Augen
\item Einsatzgebiet
\item Sensormaterialien
\item Betriebseigenschaften
\item Sondermerkmale
\item[] …
\end{anumerate}

\section{Statische Eigenschaften von Sensoren}
(keine Funktion der Zeit)\\
$\overset{x}{\longrightarrow}$Sensor$\overset{y}{\longrightarrow}$
\paragraph{(statische) Transferfunktion} (Empfindlichkeitskennlinie, Kennlinie):\\
Beziehung zwischen Eingangssignal $x$ und Ausgangssignal $y$.
$$\boxed{y=f(x)}$$
\subparagraph{Ideale Kennlinie} soll möglichst
\begin{itemize}
\item monoton
\item eindeutig
\item linear
\end{itemize}
sein (bspw. $f(x)=x$).\\
Reale Transferfunktionen weisen in der Regel nicht alle dieser Eigenschaften auf (ggf. quantisierung nötig).\\
Weiterhin ist der Messbereich (FSI: full scale input) und der Ausgabebereich (FSO) in der Regel begrenzt.
\subparagraph{Dezibel} $G[dB]=20 \log \frac{s}{s_0}$ (Einheit von FSI und FSO)\\
für $s_0$ wird dann die Normgröße für Kraft, Strom oder Spannung eingesetzt (Spannung: $1 \mu V$, Leistung: $1mV$, Schallpegel: $1\cdot 10^{-12}W/m^2$)

\subsection{Auflösung}
$\Delta x_{\text{min}} \Rightarrow$ fassbares $\Delta y$\\
Sprich: „der kleinste Schritt der Eingangsgröße liefert auch noch eine fassbare Ausgangsgröße“

\subsection{Sensitivität} (Empfindlichkeit)
$$S(x)=\frac{\diffd{y}(x)}{\diffd{x}}$$
Sensitivität abhängig von Steigung der Funktion:
\begin{itemize}
\item hohe Steigung: hohe Sensitivität
\item mittlere Steigung: mittlere Sensitivität
\item niedrige Steigung: niedrige Sensitivität
\end{itemize}
\subsection{Linearität} (Nichtlinearität)\\
Betrachtet wird die maximale Abweichung zwischen einer linearen Näherung einer nicht-linearen Funktion ($\Delta y_{\text{max}}$).\\
$\epsilon_{\text{max}}=\frac{\Delta y_{\text{max}}}{FSO}\cdot 100 \%$\\
ABB AS9

\subsection{Ansprechschwelle}
$x_{\text{min}}$\\
ABB AS6
\subsection{Totband}
ABB AS7
\subsection{Sättigung}
ABB AS8

\subsection{Wiederholgenauigkeit}
\subsubsection{Reproduzierbarkeit}
Maximales Intervall des Stimulus, bei welchem sich unter gleichen Laborbedingungen das Ausgangssignal nicht gleich ist.\\
ABB AS10\\
$\delta_R = \frac{\Delta x_{\text{max}}}{FSI}$

\subsubsection{Hysterese}
(Wenn man die Messung zwei mal von unterschiedlichen Richtungen ausführt $\to$ Abweichung in $y$-Wert abhängig von der Richtung)\\
ABB AS11

\subsection{Selektivität}
(Querempfindlichkeit)\\
Empfindlichkeit des Messwertes auf andere Größen
\begin{itemize}
\item Empfindlichkeitsänderung\\
ABB AS12
\item Nullpunktdrift\\
ABB AS13
\end{itemize}
Behebung: Kompensationssensor einsetzen.

\subsection{Rauschen}
ABB AS14

\subsection{Weitere}
\begin{itemize}
\item Langzeitstabilität
\item Lebensdauer
\item Exemplarstreuung
\item …
\end{itemize}

\section{Dynamische Eigenschaften von Sensoren}
\begin{itemize}
\item zeitlich veränderliche Messgrößen
\item diese seien adäquat (entsprechend der statischen Transferfunktion) und verzögerungsfrei weiterzugeben
\end{itemize}
Beispiel:\\
Energiespeicher
\begin{itemize}
\item Massen
\item elektrische Kapazitäten
\item Induktivitäten
\item Wärmekapazitäten
\item dissipative Energieverluste
\end{itemize}

\paragraph{Ansprechverhalten} Antwort eines Sensors auf ein zeitlich veränderliches Eingangssignal

\paragraph{Ordnung des Sensors} Anzahl vorhandener Energiespeicher (Ordnung der Differentialgleichung)
\begin{itemize}
\item 0. Ordnung: keine Zeitabhängigkeit $I(t) = U_0 \cdot \sigma (t)$
\item 1., 2.,… Ordnung: zeit- und frequenzabhängiges Ansprechverhalten
\end{itemize}

\subsection{Sprungantwort und Zeitkenngrößen}
Sprunganregung: $x(t) =\begin{cases}
0 &(t>0)\\
x_0 & (t \geq 0)
\end{cases}$\\
ABB AS15
\begin{itemize}
\item 0. Ordnung: $y(t\geq 0) = S(x_0)$
\item 1. Ordnung: (bspw. Kapazität, die auf-/entladen wird) $U_C(t) = U_0 \cdot \left\lbrace 1- \exp \left(-\frac{t}{\tau}\right)\right\rbrace$\\
$\tau = R \cdot C$ (Zeitkonstante)\\
ABB AS16\\
Halbwertszeit: $t_{0,5}$\\
Anstiegszeit (rise time): $t_r = t_{0,9}$ (also $t_{90\%}$, wann $y_{st}$ zu $90\%$ erreicht ist)\\
Einstellzeit: $t_E = 3 \tau$
\end{itemize}
\subsubsection{Dämpfung}
\begin{itemize}
\item $\delta$ Dämpfungskonstante
\item $\omega_0$ Eigenkreisfrequenz
\end{itemize}
Dämpfungsgrade:
\begin{itemize}
\item $\delta = 0$: ungedämpfte Schwingung%
\item $\frac{\delta}{\omega_0}<1$: Schwingfall
\item $\frac{\delta}{\omega_0}\geq 1$: Kriechfall (stationärer Wert wird ohne Überschwingung erreicht)
\end{itemize}
ABB AS18
\begin{itemize}
\item Totzeit $t_T$
\item Verzögerungszeit $t_V$
\item Ausgleichszeit $t_A = \tau$
\item Einstellzeit $t_E$ 
\end{itemize}
Das Toleranzband ist vorher zu bestimmen!
\begin{itemize}
\item $\delta_{opt}=0,7 \dots 0,8 \cdot \omega_0$
\item $\frac{\Delta y}{y_{st}}=\pm 4,5 \%$
\end{itemize}
\subsection{Frequenzkenngrößen}
Antwort auf ein sinusförmiges Eingangssignal: $x(t) = x_0 \cos (\omega t)$ \quad ($\omega$: Erregerfrequenz)
\begin{itemize}
\item stationärer (eingeschwungener) Zustand: $\lim_{t\to \infty} y(t) = \hat{y}_{st}(\omega) \cdot \cos (\omega t + \varphi)$\\
Nicht alle spektralen Komponenten (Frequenzen) können amplituden- und phasentreu übertragen werden.\\
\emph{stationärer Amplitudengang} (magnitude ratio): $M=\frac{\hat{y}_{st}(\omega)}{\hat{y}_{st} (\omega \to 0)}$\\
\emph{stationärer Phasengang}: $\varphi$\\
Sensoren 0. Ordnung: amplitudentreu bei $M=1$, phasentreu bei $\varphi = 0$
\end{itemize}
\subsubsection{Frequenzkenngrößen von Sensoren 1. Ordnung}
ABB AS19\\
Grenzfrequenz: $\omega_G$\\
$M=0,707$ \quad $\varphi=45^\circ$\\
Grenzfrequenz wird so angelegt, dass $\omega/\omega_G=1$
\begin{itemize}
\item $\omega/\omega_G \ll 1$:\\
$M \approx 1$\\
$\varphi \approx 0$
\item $\omega / \omega_G \gg 1$:\\
$M \to 0$\\
$\varphi \to 90^\circ$
\end{itemize}
Bandbreite: Frequenzbereich mit $M$ zwischen $0,707\dots 1,414 (\pm 3dB)$

\subsubsection{Frequenzkenngrößen von Sensoren 2. Ordnung}
ABB 20\\
$\zeta = \frac{\delta}{\omega_0}$
\begin{itemize}
\item quasistatische Anregung:\\
$\omega \ll \omega_0$: $M\approx 1$, $\varphi \approx 0$
\item Resonanz (bis $\delta = 0,71 \omega_0$):\\
\emph{Resonanzfrequenz}: $\omega = \sqrt{\omega_0^2-2\delta^2}$\\
$\varphi = -90^\circ$, $M\to M_{max}$\\
Ab Dämpfung $\zeta = 0,7$ verschwindet die Resonanz beinahe vollständig.
\item hohe Anregerfrequenz $\omega \gg \omega_0$\\
$\varphi \to -180^\circ$, $M\to 0$
\item Bandbreite: Frequenzbereich mit $M$ zwischen $0,707\dots 1,414 (\pm 3dB)$
\end{itemize}

\chapter{Temperatursensoren}
Drucksensoren
Durchflusssensoren
Binäre Positionssensoren
Flüssigkeits-Chemosensoren
…
\section{Temperaturmessverfahren}
Alle physikalischen Eigenschaften sind mehr oder weniger stark temperaturabhängig:
\begin{itemize}
\item Gasthermometer (73-543 K)
\item Flüssigkeitstherm. (203-360 K)
\item Farbstofftherm. (303-1873 K)
\item Pt100 (53-1123 K)
\item Keramik NTC, PTC (193-523 K)
\item Thermoelemente (53-2773 K)
\item Si-Halbleiter (208-573 K)
\item IR-Sensoren (223-3473 K)
\end{itemize}
Hauptsächlich 2 Methodentypen:
\begin{itemize}
\item Kontaktthermometrie (misst im Prinzip allerdings nicht die Temperatur des Mediums, sondern seine eigene $\Rightarrow$ braucht Zeit um die Temperatur bei Kontakt anzugleichen $\to$ Zeitaufwändig)
\item Pyrometrie (Analysiert Wärmestrahlung)
\end{itemize}
Die für uns wichtigsten Kontaktthermometer sind:
\begin{itemize}
\item Thermoresistive Sensoren
\item Thermoelemente
\item Halbleiterfühler verschiedener Art (Sperrstrom oder Schwellspannung)
\end{itemize}
\section{Zeitverhalten eines Widerstandsensors}
ABB AS21\\
(Verhalten wie Sensor 1. Ordnung)
\begin{itemize}
\item Verzögerungszeit $t_V$
\item Ausgleichszeit $t_A = \tau$
\item Halbwertszeit $t_{50}$
\item Anstiegszeit $t_{90}$
\item Einstellzeit $t_E=3 \tau$
\end{itemize}
Luft: sehr langsame Anpassung der Temperatur am Sensor\\
Flüssigkeit: schnellere Anpassung (ca. 10 mal schneller als Luft)\\
Festkörper: etwas langsamere Anpassung

\section{Thermoresistive Sensoren}
ABB AS 22\\
Temperaturkoeffizient des Widerstands:
$$TK_R=\frac{1}{R}\cdot \frac{\diffd{R}}{\diffd{T}}$$
Unterscheidung der Sensoren:
\begin{itemize}
\item Kaltleiter $TK>0$ (PTC: positive temperature coefficient)\\
bspw. Metallwiderstände ($Ni$ und $Pt$)\\
Vorteil: in kleinen Temperaturbereichen ist Anstieg linear
\item Heißleiter $TK<0$ (NTC: negative temperature coefficient)
\end{itemize}

\subsection{Metallwiderstände}
Je nach Metall (Platin, Kupfer, Nickel) ist die Auslenkung unterschiedlich. Platin hat eine optimale Auslenkung, ist aber teuer. Etwas schlechter, aber billiger ist Kupfer. Nickel als alternative hat ab der Curie-Temperatur ($352 ^\circ C$) eine sehr schlechte Auslenkung.\\
Die Metalle sollten folgende Eigenschaften haben:
\begin{itemize}
\item korrosionsfest
\item alterungsfest
\item hohe $T_S$ (Schmelztemperatur)
\item hoher $\varrho$ (elektrischer Widerstand)\\
(da kleine Änderungen einer kleinen Größe gemessen werden soll)
\item duktil
\end{itemize}
\begin{tabular}{l l l}
Metall & $\varrho $ & TK\\
&$ 10^{-8}\Omega \mathrm{m}$&$10^{-3}K^{-1}$\\
\hline
Platin (Pt) & $10,6$ & $3,92$\\
Kupfer (Cu) & $1,67$ & $4,33$\\
Nickel (Ni) & $6,84$ & $6,75$
\end{tabular}\\
Man sieht: bei Metallen ist der Temperaturkoeffizient sehr klein.\\
Sensoren:
\begin{itemize}
\item Keramikkörper, in den Platinwendel (Drähte) lose in Löcher eingelegt werden
\item Keramikdorn, um den Platinwendel (Drähte) umgewickelt werden (abgedeckt mit Glasrohr)
\item alternative (zu den Draht-Lösungen): im Sensor wird eine Platindünnschicht auf Keramikträger aufgetragen (sehr wenig Platin $\to$ deutlich günstiger; Nachteil: dünne Platinschicht liegt dicht auf Keramik. Bei Temperaturänderungen dehnt sich Keramik und Platin unterschiedlich schnell, das führt zu mechanischen Spannung)
\end{itemize}

\subsubsection{Pt100}
gebräuchlichster Temperatursensor
\begin{itemize}
\item Spezifikation: DIN EN 60571
\item Temperaturbereich (Dauerbetrieb): $-7^\circ C$ bis $+550^\circ C$
\item Temperaturkoeffizient: $Tk=3850 \;\mathrm{ppm/K}$
\item Langzeitstabilität: Max. $R_0$-Drift: $0,04\%$ nach $1000h$ bei $500^\circ C$
\item Erschütterungsfestigkeit: Min. $40g$ Beschleunigung bei $10$ bis $2000\;\mathrm{Hz}$
\item Umgebungsbedingungen: Ungeschützt nur im trockener Umgebung einsetzbar
\item Isolationswiderstand: $>10 M\Omega$ bei $20^\circ C$, $>1 M\Omega$ bei $500^\circ C$
\item Selbsterwärumng: $0,2K/mW$ bei $0^\circ C$
\item Ansprechzeit: \\
Bewegtes Wasser ($v=0,4m/s$): $t_{0,5}=0,3s; t_{0,9}=0,8s$\\
Lufstrom ($v=1m/s$): $t_{0,5}=3,0s;t_{0,9}=9,0s$
\item Messstrom: $100 \Omega$: $0,1$ bis $0,4 mA$
\end{itemize}
Vorteile:
\begin{itemize}[label=$+$]
\item für Präzisionsmessungen geeignet ($\pm 0,3\%$)
\item sehr gute Langzeitstabilität
\item gute Linearität (bis $150^\circ C$)
\item PT-Sensoren: $-270^\circ C$ bis $+850^\circ C$\\
mech. sehr stabil, Dünnfilmsensoren
\item Ni-Sensoren: kostengünstige Dickschichtsensoren
\end{itemize}
Nachteile:
\begin{itemize}[label=$-$]
\item kostenaufwendig
\item wegen Gehäuse meist langsam
\item Spannungsversorgung erforderlich
\item Ni-Sensoren: $-60^\circ C$ bis $+180^\circ C$
\item keine punktförmige Messung
\item Selbsterwärmung
\end{itemize}
\subsubsection*{Schalttechniken}
\begin{itemize}
\item Zweileitertechnik:
Auslesen des Spannungsabfalls über Pt100-Sensor direkt\\
Problem: Widerstände der Drähte von/zum Pt100-Sensor können Signal verfälschen.
\item Vierleitertechnik:\\
Auslesen des Spannungsabfalls mit einer parallelen Leitung zusätzlich zur Versorgungsleitung des Pt100. Damit ist die Messung Stromfrei. \\
Nachteil: 4 Leiter mit Widerständen, die verfälchen können.\\
Lösung: Dreileitertechnik: nicht komplett parallele Spannungsmessung, sodass ein Widerstand weg fällt.
\end{itemize}
\subsection{Kalteiter (PTC)}
\subsubsection{Thermositoren} (temperaturabhängige Widerstände)
\begin{itemize}[label=$+$]
\item preiswert
\item großer Messwiderstand
\item $-30^\circ C$ bis $+350^\circ C$ ($\pm 5^\circ C$)
\item große Empfindlichkeit
\item ohne Gehäuse realisierbar $\Rightarrow$ kurze Ansprechzeit
\end{itemize}
\begin{itemize}[label=$-$]
\item stark nichtlineare Kennlinie
\item ohne Gehäuse mechanisch anfällig
\item Spannungsversorgung erforderlich
\item keine punktförmige Messung
\end{itemize}
\subsubsection{schaltende Sensoren}
\begin{itemize}
\item Überstrom-Überwachung
\item Übertemperatur-Überwachung
\item Niveaufühler für Flüssigkeiten
\end{itemize}
\subsubsection{Si-Temperatursensoren}
Widerstand steigt mit Temperatur $\to$ Kaltleiter.
\begin{itemize}[label=$+$]
\item preiswert
\item großer Messwiderstand
\item große Empfindlichkeit
\end{itemize}
\begin{itemize}[label=$-$]
\item keine lineare Kennlinie
\item DMS-Effekt (Dehnungs-Mess-Streifen)
\item Spannung erforderlich
\item Gehäuse $\Rightarrow$ langsam
\item keine punktförmige Messung
\end{itemize}

\subsubsection{Spreading-Resistance-Sensor}
$R=\frac{\varrho(T)}{\pi \cdot d}$\\
ABB AS23 (Ersatzschaldbild)

\subsubsection{Temperatur IC}
Diodengleichung: $I=I_0\cdot \left[ \exp\left(\frac{e \cdot U}{k \cdot T}\right) -1 \right]$ mit $I_0$: Sperrstrom

\subsection{Heißleiter (NTC)}
\begin{itemize}[label=$+$]
\item breites Spektrum an $R$
\item großer TK ($3 \dots 6 \% /K$)
\item[$-$] $-50^\circ C \dots +100^\circ C (\pm 5^\circ C)$
\end{itemize}
\begin{itemize}
\item Thermistoren: Bestehen aus Keramiken (oxyde).
\item Temperaturfühler
\item Zeitverzögerungsglied
\item Spannungsstabilisierung
\end{itemize}

\subsubsection{Thermoelemente}
ABB AS24\\
Beispiel: Draht führt durch zu Eis/Kerze (Kupfer), dazwischen ist dann Konstantan. Der Spannungsunterschied wird gemessen.\\
Seebeck-Effekt: $U_{th}=\alpha_{th} \cdot \Delta T$ mit $U_{th}$: Thermospannung und $\alpha_{th}$: Thermokraft

\subsubsection{Genormte Thermoelemente}
\begin{tabular}{L{0.2} | l L{0.1} | l | L{0.3}}
Thermopaar & Kurzzeichen & lin. Thermosp., $\mu V/K$ & T-Bereich ($^\circ C$) & Eigenschaften\\ \hline 
Eisen-Konstantan (Fe-CuNi) & Typ J (Typ L) & 53 & -250 bis 700 (900) & billig, rostet, empfindlich gegen schwefel- und kohlenwasserstoffhaltige Luft\\\hline
Kupfer-Konstantan (Cu-CuNi) & Typ T (Typ U) & 443 & -250 bis 300 & rostet nicht, Cu-CuNi verzundert\\\hline
Nichkelchrom-Nickel (NiCr-Ni) & Typ K (Typ N) & 41 & -200 bis 1000 (1300) & hohe Thermokraft, recht linear, empfindlich gegen S-haltiges Gas\\\hline
Platinrhodium-Platinum (PtRh-Pt; Pt/Rh=87/13) & Typ R & 6,5 & -50 bis 1500 & bei Rotglut einsetzbar, teuer, oxidationsbeständig\\\hline
Platinrhodium-Platinum (PtRh-Pt; Pt/Rh=90/10) & Typ S & 6,5 & -50 bis 1500 & wie Typ R (angesächsisch), aber mehr im dt. Raum\\\hline
PtRh-Schenkel mit verschiedenem Rh-Gehalt (70/30 - 94/6) & Typ B & 6,3 & 0 bis 1600 (1800) & unempfindlich gegen Verunreinigungen, T der Vergleichsstelle muss nicht konstant sein\\\hline
WRh-Schenkel mit verschiedenem Rh-Gehalt (75/25 - 97/3) & Typ G & 15 & 0 bis 2200 & T der Vergleichsstelle muss nicht konstant sein, nur in neutraler und reduzierender Atmosphäre
\end{tabular}

\subsubsection{Aufbau von Thermoelementen}
\begin{itemize}
\item Kabelthermoelemente
\begin{itemize}[label=$+$]
\item sehr schnell
\item[$-$] ungeschützt
\end{itemize}
\item Mantelthermoelemente (im Prinzip Kabel inkl. Mantel) $\to$ etwas langsamer, aber geschützter
\item Isolierte Messstelle (TI)
\begin{itemize}
\item[$+$] elektrische und galvanische Trennung
\item[$-$] lange Zeitkonstante
\end{itemize}
\item mit Mantel verbundene Messstelle (TM) [an Messstelle liegt Thermoelement „blank“]
\begin{itemize}
\item[$+$] kurze Zeitkonstante
\item[$-$] galvanische Verbindung
\end{itemize}
\end{itemize}
\subsubsection{Ansprechzeiten, Vorteile, Nachteile der Thermoelemente}
\begin{itemize}[label=$+$]
\item sehr großer T-Bereich
\item kleine Gesamtabmessungen
\item große Flexibilität
\item einfacher, sehr robuster Aufbau
\item unempfindlicher gegen Hitze, Druck und Vibration
\item kleinste Ansprechzeiten
\item praktisch punktförmige Messstellen
\end{itemize}
\begin{itemize}[label=$-$]
\item reduzierte Genauigkeit
\item sehr kleine Thermospannung
\item Temperatur der Anschlussstelle muss bekannt sein
\end{itemize}

\subsection{Bimetallthermometer}
Zwei unterschiedliche Metalle sind in einem Streifen miteinander verbunden. Dadurch, dass sie sich bei Temperatur unterschiedlich verformen, kann durch die Auslenkung die Temperatur berechnet werden.\\
$f$: Auslenkung der Spitze des Streifens,\\
$s$: Dicke des Streifens (beider Metalle),\\
$L$: Länge des Streifens\\
$\delta$: spezifische Ausbiegung
$$f=\frac{L^2}{s}\cdot \delta \cdot \Delta T$$
Einsatzbeispiel: Schalter wird durch den Bimetallstreifen betätigt, wenn die Temperatur bewirkt, dass er sich verformt: schaltende Sensoren (Überlastungsschutz, Temperaturregelung).

\section{Leitlinien für die Temperaturmessung}
\begin{itemize}
\item Messort
\item Gehäuseeinfluss
\item Ansprechzeit
\item Eigenerwärmung
\item Absolute Genauigkeit
\end{itemize}
Für hohe Ansprüche an die Auflösung und Genauigkeit: \fbox{Pt-Fühler}\\
Für kleinste Messstellen und schnelles Ansprechen: \fbox{Thermoelemente}
\paragraph{Messprobleme bei:}
\begin{itemize}
\item Dickem Gehäuse
\item Große Entfernung zwischen Objekt und Messstelle
\item Objekte geringer Masse (mit massiven Fühlern $\to$ Verfälschung der Messung durch Angleichung der zu messenden Temperatur an die Fühlertemperatur)
\item T-Messung bei Luftströmungen
\end{itemize}

\chapter{Druckmessung}
\section{Druck als Messgröße}
Normaldruck: $1013\;\mathrm{hPa}$\\
$\text{Druck}=\frac{\text{Kraft}}{\text{Fläche}} \quad [p]=1\;\mathrm{Pa}=\frac{1 \;\mathrm{N}}{1 \;\mathrm{m^2}} \quad 1 \;\mathrm{bar} = 10^5 \;\mathrm{Pa}$\\
\begin{tabular}{l | l}
Druckbereich & Anwendung\\\hline
$<40 \unit{mbar}$ & Füllstand in Waschmaschine, Geschirrspüler\\
$100\unit{mbar}$ & Staubsauger, Filterüberwachung, Durchflussmessung\\
$200\unit{mbar}$ & Blutdruckmessung\\
$1\unit{bar}$ & Barometer, Korrektur für Zündung und Einspritzung\\
$0\dots 1 \unit{bar}$ & Ansaugunterdruckmessung\\
$10\unit{bar}$ & Öldruck und Luftdruck für Bremsen, Kühlmaschinen\\
$50 \unit{bar}$ & Pneumatik, Industrieroboter\\
$500 \unit{bar}$ & Hydraulik, Baumaschinen\\
$3000 \unit{bar}$ & Hochdruckschneiden mit Wasser
\end{tabular}
\section{Arten der Druckmessung}
Absolutdruck: Ausgehend von keinem Druck ($0$)\\
Unterdruck/Überdruck: Ausgehend vom Relativdruck (festgelegt: $p_{atm}$, bspw. Atmosphärendruck)\\
Differenzdruck: Druck zwischen zwei Messungen $p_{diff}=p_1-p_2$
\begin{itemize}
\item Absolut-Drucksensor:\\
zu messender Druck auf einer Membran gegenüber dem Vakuum
\item Relativ-Drucksensor:\\
zu messender Druck auf eine Membran gegenüber dem Umgebungsdruck
\item Differenz-Drucksensor:\\
die beiden Drücke auf beiden Seiten der Membran
\end{itemize}
\section{Prinzipien der Druckmessung}
\begin{itemize}
\item Verformung unter Druck: magnetoelastische und Membransensoren
\item Druckabhängigkeit physikalischer Eigenschaften: piezoelektrische Kristallsensoren
\end{itemize}
Störeinflüsse:
\begin{itemize}
\item mechanische Spannungen
\item Temperatureffekte
\end{itemize}
\subsection{Membransensoren}
Membran biegt sich unter Druck $p$ (mit einer Durchbiegung $\Delta x$)

\subsubsection{Druckschalter}
unter Druck wird Schalter betätigt. Abhängig von den eingebauten Membranen und Druckfedern.
\begin{itemize}
\item[$+$] Messbereich: $\unit{mbar}\dots\unit{kbar}$
\item[$-$] relativ große Hysterese (Einschalt- und Ausschaltdruck unterscheiden sich)
\end{itemize}
\subsubsection{Kapazitiver Drucksensor}
Auf Membran ist Messelektrode angebracht. Durch die Verformung der Membran variiert der Abstand zwischen der Messelektrode und der Referenzelektrode. Dadurch verändert sich die Spannung.\\
$C=\varepsilon\varepsilon_0\frac{A}{d}$\\
Typische Abmessungen:
\begin{itemize}
\item Äußerer Druckmesser: $32,4\unit{mm}$
\item Dicke des Grundkörpers: $5\unit{mm}$
\item Dicke der Membran (je nach Messbereich): $0,2\dots 2,8 \unit{mm}$
\item Abstand zwischen Membran und Grundkörper: $40\unit{\mu m}$
\item Durchbiegung bei Nenndruck: max. $8 \unit{\mu m}$
\end{itemize}
Eigenschaften:
\begin{itemize}[label=$+$]
\item sehr empfindlich
\item geringer Energieverbrauch (da nicht generierender, sondern modulierender Sensor)
\item geringe T-Abhängigkeit
\item sehr einfach, robust
\item nicht empfindlich für Überdruck
\item besonders für niedrige $\unit{P}$: $10^{-7}\dots 1 \unit{bar}$
\item gasartunabhängig
\item Genauigkeit: $0,2\%\text{ (Ker)}\dots2\%\text{ (Me)}$ vom Messwert
\item lineare Kennlinie
\end{itemize}
\begin{itemize}[label=$-$]
\item Kapazitätsänderung: Bruchteile von $\unit{pF}$
\item Streukapazitäten-, staub- und feuchteempfindlich
\item relativ komplexe Elektronik
\end{itemize}

\subsubsection{Kapazitive MEMS-Sensoren}
Mikro-Elektro-Mechanische Systeme: Ähnlich wie kapazitive Drucksensoren, bloß deutlich kleiner.
\begin{itemize}
\item Drucksensor
\item Beschleunigungssensor
\end{itemize}

\subsubsection{Dehnungsmessstreifen (DMS)}
Messstreifen mit einer Länge $l$ und Durchschnitt $A$:\\
$R=\rho \frac{l}{A}$\\
Da sich der Durchschnitt kaum ändert: $\frac{\Delta R}{R}=\tred{k}\cdot \frac{\Delta l}{l}$\\
Der \tred{k}-Faktor:\\
\begin{tabular}{l | c}
Werkstoff & K-Faktor\\\hline
Konstantan & 2,0\\
Nichrom & 2,2\\
FeNi & 2,5\\
„Iso-Elastic“ & 3,6\\
PtW & 4,0
\end{tabular}\\
Der k-Faktor ist sehr stark von der Temperatur abhängig!

\subsubsection*{Membran-DMS-Sensorelement}
Widerstand-Anordnung an der Membran ist zu vergleichen mit einer Wheatstone-Brücke (mit $R_1$ und $R_2$ parallel zu $R_3$ und $R_4$), wobei $R_1=R_2=R_3=R_4$.\\
$R$: $120\dots5000\unit{\Omega}$\\
$\Delta R \sim 0,5\dots 5\%$
\begin{itemize}
\item Metall-Dünnfilm-DMS
\item polySi-Membran-DMS
\end{itemize}
\begin{itemize}[label=$+$]
\item kleine Abmessungen möglich
\item für statische und dynamische Messungen geeignet
\end{itemize}

\subsubsection{Piezoresistive Metall-Mebran-Drucksensoren}
\begin{itemize}[label=$+$]
\item höchste Genauigkeit: $0,2 \ldots 0,3\%$, sehr kleine Hysterese $<0,05\%\mathrm{FSO}$
\item hervorragende Temperaturstabilität
\item überdrucksicher
\item sehr hohe Medienverträglichkeit
\item robuste Gehäuseausführungen möglich
\item direkter Kontakt mit Druckmedium
\item gute Langzeitstabilität
\end{itemize}
\begin{itemize}[label=$-$]
\item für Drücke $<5\mathrm{bar}$ nicht geeignet
\item für $T > 140^\circ C$ nicht geeignet
\item aufwendige Fertigung
\item sehr teuer (hoher Fertigungsaufwand)
\item geringe elektrische Durchschlagsfestigkeit
\item anfällig für elektro-magnetische Störeinkopplung
\item Miniaturisierung eingeschränkt
\end{itemize}
Anwendung für Präzisions-Druckmessungen insbesondere in Hochdruckbereichen:
\begin{itemize}
\item Pneumatik und Hydraulik
\item Mess- und Regeltechnik
\item Verfahrenstechnik der Chemie, Pharmaindustrie, …
\end{itemize}

\subsubsection{Silizium-Membran-Drucksensoren}
Si-Messmembran mit poly-Si-Dünnfilm DMS
\begin{itemize}[label=$+$]
\item Ideal für niedrige Nenndruckbereiche
\item gutes Preis/Leistungsverhältnis
\item vergleichsweise kleine Abmessungen
\item Integration mit der Signalverarbeitungselektronik
\item hoher k-Faktor
\item keine Kriecheffekte
\item keine Isolierschicht
\item hohe Ausgangssignale
\item hohe Eigenfrequenz
\end{itemize}
\begin{itemize}[label=$-$]
\item hohe Temperatur-Abhängigkeit
\item hoher Aufwand für Linearisierung und Temperatur-Kompensation
\item empfindlich gegen äußere Einflüsse
\item geringere Genauigkeit als Metall-DMS
\item Einsatz nur bis max. $125^\circ C$
\item Hysterese $0,15\% \mathrm{FSO}$
\end{itemize}
Anwendungsbereiche:
\begin{itemize}
\item keine hohen Genauigkeitsansprüche
\item insbesondere im Niedrigdruckbereich
\item Haushalt (keine Temperatur-Schwankungen)
\end{itemize}

\subsection{Piezoelektrische Effekte}
$q=d\cdot F$ mit $d$ … piezoelektrischer Koeffizient\\
Prinzip: Atom-Anordnung ist anders bei unterschiedlichem Druck
\begin{itemize}
\item normal: sehr symmetrische Anordnung
\item Longitudinal-Effekt (Druck von oben und unten, dort wo Ladung angelegt wird): Atome werden zusammengedrückt, so sind teilweise oberhalb unterschiedliche Ladungen (positiv/negativ)
\item Transversal-Effekt (Druck von oben und unten, Ladung liegt seitlich an): Atome werden zusammengedrückt und gedreht, dadurch entstehen wieder unterschiedliche Ladungen
\item Schub-Effekt (Ladungsplatten (oben/unten) werden seitlich verschoben): Atome verschieben sich
\end{itemize}
Spannung über dem Sensorelement: $U=\frac{d}{\varepsilon_r}\cdot \frac{F \cdot l}{\varepsilon_0 \cdot A}$

\subsubsection*{Piezoelektrische Materialien}
\begin{tabular}{L{0.3} | L{0.2} | L{0.15} | L{0.15} | L{0.2}}
Material & Dielektrizitätskonstante $\varepsilon_{33}$ & Piezoelktrische Dehnungskonstante $d_{31}[\unit{pC/N}]$ & Curie-Temperatur $[\unit{^\circ C}]$ & typische Einsatz-Temperatur\\\hline
Quarz & $4$ & $2,3$ & $560$ & $-200\ldots 300\unit{^\circ C}$\\
\multicolumn{5}{l}{$\to$ exzellente Linearität und Langzeitstabilität, keine Hysterese, keine Pyroelektrizität.}\\ \multicolumn{5}{l}{Aber: quarz ist polymorph und stärker temperaturabhängig}\\\hline
Blei-Zirkonat-Titanat Keramik (PZT) P8 & $1600$ & $300$ & $370$ &\\\hline
Bleimetaniobat & $200$ & $75$ & $540$&\\\hline
Turmalin & $7$ & $1,8$ & $T_m\geq 900$ & bis $650\unit{^\circ C}$\\\hline
Keramik P14 & $160$ & $20$ & $680$ &\\\hline
Lithiumniobat & $85$ & $80$ & $1200$ & bis $700\unit{^\circ C}$\\\hline
$BaTiO_3$-Keramik & $2000-5000$ & $190$ & $120$ &
\end{tabular}

\subsubsection*{Piezoelektrische Drucksensoren}
\begin{itemize}[label=$+$]
\item nahezu weglose Kraftaufnehmer (nur wenige $\unit{\mu m}$)
\item hoher Wirkungsgrad (mechanische in elektrische Energie)
\item einfacher mechanischer Aufbau: flach, klein, leicht
\item sehr großer Messbereich ( $0 \ldots 100 \unit{MPa}$
\item hohe Linearität des piezoelektrischen Effekts
\item Auflösung: $< 10^{-7} \unit{FSI}$, Genauigkeit: $1\%$
\item großer Frequenzbereich $f > 100 \unit{kHz}$
\item kleine Bauweise
\item Betriebstemperaturbereich: $-196 \ldots + 240 \unit{^\circ C}$
\end{itemize}
\begin{itemize}[label=$-$]
\item Kriechströme
\item statische Messungen nur beschränkt möglich
\end{itemize}
Anwendungsbereich:
\begin{itemize}
\item insbesondere Messung hoher Drücke bei erhöhter Temperatur
\item Werkzeugmaschinen
\item Kunststoffverarbeitende Industrie
\item KFZ (Brennkammerdrucksensoren)
\item Sport (zur mehrdimensionalen Messungen)
\item Ballistik
\item Beschleunigungsmessung
\end{itemize}

\subsection{Beschleunigungssensoren}
Prinzip: Trägheitskraft auf seismische Masse\\
Wirkungsweise:
\begin{itemize}
\item kapazitive Änderung (durch Verschiebung) oder
\item piezoelektrisch:
\begin{itemize}
\item Aufnehmer vom Kompressionstyp
\item Aufnehmer nach dem Scherkonstruktionsprinzip
\item Aufnehmerelement unter Biegebeanspruchung
\end{itemize}
\end{itemize}
\subsubsection{Airbag-Beschleunigungssensor}
(auf Basis der Scherkonstruktion)
\begin{itemize}
\item Hohe Empfindlichkeit: bis $1 \unit{nC/g}$
\item Betriebstemperaturbereich: $-55 \dots 120\unit{^\circ C}$
\item Empfindlichkeitsbereich: bis $\pm 500\unit{g}$
\item Frequenzbereich: ab $1 \unit{Hz}$ bis $0,3 f_R$
\end{itemize}
\subsubsection{Bimorphe Beschleunigungssensoren}
(auf Basis der Biegebeanspruchung) Die seismische Masse ist das Piezoelement selbst.
\begin{itemize}
\item Masse: $4,5 \unit{g}$
\item Empfindlichkeit: $1\unit{V/g}$
\item leicht, klein, hochauflösend, kostengünstig
\end{itemize}

\subsubsection{Anwendungen}
\begin{itemize}
\item Generelle Maschinenschwingungen
\item Verpackungsindustrie
\item Radialbeschleunigungssensoren (Eisenbahn)
\item Crashtests
\item Airbags im Auto
\item Karosserieanalyse
\item Sport (Aktivitätssensor)
\item Lagesensor
\end{itemize}

\subsection{Magnetoelastische Sensoren (Induktivitätssensoren)}
Prinzip: Spuleninduktivität (Villari-Effekt, Differentialtransformator)\\
$L=\mu_0 \mu_r \frac{N^2 A}{l}$

\begin{itemize}
\item statische und quasi-statische Messungen
\item Messeffekte sind sehr groß (keine große Verstärkung nötig)
\item keine Rückwirkung auf Messobjekte
\item Einsatzbereich: $1 \unit{kN}\ldots 2,5 \unit{MN}$
\item Messfehler: $1\%$
\end{itemize}

\chapter{Durchflusssensoren}
Rohrströmung
\begin{itemize}
\item laminare Rohrströmung\\
Strömung ohne Turbelenz
\item turbulente Rohrströmung
\end{itemize}
Reynolds-Zahl: $\mathrm{Re}=\frac{\nu \cdot D \cdot \rho}{\eta}$\\
Fließgeschwindigkeit im Rohr ist je nach laminarer oder torbulenter Strömung unterschiedlich verteilt. Bei der torbulenten Strömung ist das Profil (über den Querschnitt des Rohrs) ausgeglichener, bei laminarer zentriert (mit größerem Ausschlag).

$\Rightarrow$ eine Ein-Punkt-Messung gibt keine Aussage über die gesamte Strömung.

\section{Prinzipien der Durchflussmessung}
\begin{itemize}
\item Volumenstrom\\
$q_v= \frac{\Delta V}{\Delta t}$
\begin{itemize}
\item Volumenzähler
\begin{itemize}
\item Volumetrische Zähler $\Delta V$
\begin{itemize}
\item Verdrängungszähler
\item Auslaufzähler
\end{itemize}
\item Strömungszähler $\Delta V = A \cdot \nu$
\begin{itemize}
\item Messflügel
\item Drosselgeräte (Wirkdruck $\Delta P$)
\item Verschiedene physikalische Effekte
\end{itemize}
\end{itemize}
\end{itemize}
\item Massenstrom\\
$q_m = \frac{\Delta m}{\Delta t}$\\
Aussagekräftig für die Frage der Menge die durchgeflossen ist.
\begin{itemize}
\item Massendurchflussmessung\\
Coriolis, thermisch
\end{itemize}
\end{itemize}

\section{Auslaufzähler}
\begin{itemize}
\item Kippzähler: Zählung der Kippvorgänge\\
Flüssigkeit fließt in eine von zwei Messkammern. Wenn Messkammer gefüllt ist, schwenkt die Messkammer zur Seite (und entleert sich), dadurch fließt Flüssigkeit nun in die andere Messkammer. Es wird die Anzahl der Kippvorgänge gezählt.
\item Trommelzähler: Zählung der Umdrehungen\\
(Vergleich Wassermühle mit Eimern statt Schaufel)
\end{itemize}
Eigenschaften:
\begin{itemize}
\item Fluideigenschaften: nur Flüssigkeiten
\item Messgenauigkeit: $\pm 0,1 \ldots 0,5 \%$ v.S. (vom Sollwert)
\item Fehlerquellen: Ungemessene Flüssigkeitsmengen während des Kippvorgangs\\
$\to$ negative Messabweichung
\end{itemize}
\begin{itemize}[label=$+$]
\item hohe Genauigkeit
\item Messung stark verschmutzter Flüssigkeiten möglich
\item niedriger Grenzwert (kleinste Mengen können verarbeitet werden)
\item geringe Druckverluste
\end{itemize}
\begin{itemize}[label=$-$]
\item Flüssigkeiten mit großer Viskosität des Mediums nicht möglich
\item nicht für Gase geeignet
\item erzeugen pulsierende Strömung
\end{itemize}
\section{Verdrängungszähler}
Volumenzähler mit beweglichen Kammerwänden
\begin{itemize}
\item Ovalradzähler
\item Drehkolben-Gaszähler
\item Zahnradzähler
\item Zahnringzähler
\item Ringkolbenzähler
\end{itemize}
Es wird gezählt, wie viel sich der Zähler bewegt/verdreht.\\
Eigenschaften:
\begin{itemize}
\item Fluideigenschaften: Flüssigkeiten und Gase
\item Messgenauigkeiten: $\pm 0,5\%$ v.S.
\item Fehlerquellen: hauptsächlich durch Spaltverluste (Verschleiß)
\end{itemize}
\begin{itemize}[label=$+$]
\item großer Messebereich $0,15$ bis $150 \unit{l/s}$
\item Betriebsdruck bis $30 \unit{bar}$
\item hohe Genauigkeit über großen Messbereich
\item weitgehende Viskositätsunabhängigkeit
\item lange Lebensdauer (bei sorgfältiger Filterung des Messgutes)
\item keine Störeinflüsse über die Strömungsprofile
\end{itemize}
\begin{itemize}[label=$-$]
\item extrem empfindlich gegenüber Verschmutzungen\\
$\to$ sorgfältige Filterung des Messgutes notwendig
\item pulsierende Strömung
\item Druckverluste bis $1 \unit{bar}$
\item Verschleiß
\item Blockieren der Leitung im Falle einer Verstopfung oder Verklemmung
\item Erhebliche Geräuschentwicklung
\end{itemize}
\section{Volumenzähler mit Messflügeln}
Drehzahl des Turbinenrades: $f \sim \nu \sim q_v$
\begin{itemize}
\item Turbinenradzähler axial/senkrecht
\item Flügelradzähler
\end{itemize}
Eigenschaften:
\begin{itemize}
\item Fluideigenschaften: reine Flüssigkeiten und Gase
\item Messgenauigkeit: $\pm 2-5\%$ v.S.
\item Fehlerquellen:
\begin{itemize}
\item durch Lagerreibung, Strömungsprofil, von der Viskosität abhängige hydraulische Verluste
\item Verwirbelung, Strömungspulsation und Druckschwankungen
\end{itemize}
\end{itemize}
\begin{itemize}[label=$+$]
\item hohe Genauigkeit bei definierten Viskositäten ($2-5\%$)
\item Temperaturen bis $350\unit{^\circ C}$
\item für große Durchflussmengen geeignet
\item Messungen bei hohem Druck bis $>600 \unit{bar}$ möglich
\item längere Lebensdauer der Lager bei senkrechten Bauarten
\end{itemize}
\begin{itemize}[label=$-$]
\item Viskosität muss bekannt sein
\item Beruhigungsstrecke vor und nach der Turbine erforderlich
\item empfindlich gegenüber Verschmutzungen $\to$ Schmutzfilter
\item aufgrund von Lagerreibung und Massenträgheit der Rotoren können Schleichmengen nicht erfasst werden
\item Strömungsprofil-abhängig
\end{itemize}
\section{Schwebekörper-Durchflussmesser}
Prinzip: Schwebekörper in einen Röhre. Von unten fließt das Medium ein. Die Steighöhe des Schwebekörpers $h$ ist ein Maß für die Strömungsgeschwindigkeit. Auf sie wirken 3 Kräfte:
\begin{itemize}
\item Gewichtskraft $\downarrow$: $F_G=m\cdot g$
\item Auftriebskraft $\uparrow$: $F_A=V\cdot \rho$
\item Kraft des Strömungswiderstandes $\uparrow$: $F_S\sim \frac{q_V}{h}$
\end{itemize}
Eigenschaften:
\begin{itemize}
\item Fluideigenschaften: Flüssigkeiten und Gase
\item Messgenauigkeit: $\pm 2\%$ v.E. (von Endwert)
\item Fehlerquellen: Abhängigkeit von der Schwerkraft, empfindlich auf Pulsation und Vibration
\end{itemize}
\begin{itemize}[label=$+$]
\item einfacher, leicht kontrollierbarer Aufbau
\item keine Vorlaufstrecke (Beruhigungsstrecke) erforderlich
\item robust, einfache Montage, leicht einstellbare Grenzwerte
\item günstiges Preis/Leistungsverhältnis
\item auch für aggressive Medien geeignet
\item näherungsweise lineare Skala
\item geringer Druckverlust
\item leichte elektronische Positionserfassung (z.B. induktiv)
\item besonders für kleine Durchflussmengen
\end{itemize}
\begin{itemize}[label=$-$]
\item kann empfindlich gegen Verschmutzung sein
\item senkrechte Einbaulage erforderlich
\end{itemize}

\section{Differenzdruckverfahren (Wirkdruckverfahren)}
Rohr mit Verengung. Differenz zwischen den Drücken vor und nach der Blende wird gemessen (Wirkdruck):
$$q_v\sim \sqrt{p_1-p_2}$$
Varianten:
\begin{itemize}
\item Blende
\item Düse (leicht führende Blende)
\item Venturidüse
\end{itemize}
Eigenschaften:
\begin{itemize}
\item Fluideigenschaften: für reine Flüssigkeiten und Gase (keine Mehrphasenströmungen)
\item Messgenauigkeit: $<1\%$ v.E.
\item Fehlerqullen: Fehler durch Verschleiß der Blenden, empfindlich für Strömungsprofil und Verschmutzung des Messmediums
\end{itemize}
\begin{itemize}[label=$+$]
\item einfach, robust und wartungsarm
\item international genormt, keine Kalibrierung erforderlich
\item besonders für große Durchflussmengen geeignet
\item unter extremen Bedingungen (P,T, Überlastbetrieb, Medium) einsetzbar
\item äußerst langzeitbeständig (je nach Aggressivität des Mediums)
\item extrem große Nennweitenpalette bei nur minimaler Preisänderung
\end{itemize}
\begin{itemize}[label=$-$]
\item nur in voll gefüllten Rohren ($D > 50\unit{mm}$) anwendbar
\item nichtlinearer Effekt: Wirkdruck steigt quadratisch mit Durchfluss
\item hoher Druckverlust (Blende $35\%-90\%$, Düse $6\%-25\%$ des Wirkdrucks)
\item große Ein- und Auslaufstrecken erforderlich
\end{itemize}

\section{Staudruckverfahren}
Es wird der Druck auf einen Staukörper gemessen (zwei Röhrchen in dem „Strahl“, zwischen denen die Druckdifferenz gemessen wird)\\
$v=\sqrt{\frac{2}{\rho}\cdot \Delta p}$ \qquad $\Delta p$: unterschied des Drucks\\
$q_V = A \cdot \nu$\\
Eigenschaften
\begin{itemize}
\item Fluideigenschaften: vor allem Gase
\item Messgenaugikeit: $(\pm 3) \pm 1 \%$ v.S.
\end{itemize}
\begin{itemize}[label=$+$]
\item robust und wartungsarm
\item punktförmige Messung der Strömungsgeschwindigkeit
\item geringe Druckverluste
\end{itemize}

\begin{itemize}[label=$-$]
\item nichtlineare Abhängigkeit der Strömungsgeschwindigkeit vom Druck\\
$\to$ hohe Messfehler im unteren Bereich
\item Gesamtdurchfluss ergibt sich durch das Abtasten der Geschwindigkeitsverteilung über Querschnitt und anschließende Integration
\end{itemize}

\section{Wirbelzähler}
Zählt Verwirbelungen, die sich hinter einem Objekt (Störkörper) in einem Strom (in Rohr) ergeben (Vergleich: Flagge an Mast).\\
Wirbelablösefrequenz: $f \sim \frac{v}{d}$\\
Messverfahren:
\begin{itemize}
\item Drucksensoren an Rohrwandung
\item Drucksensoren an Störkörper
\item Vorgeheizte Thermistatoren an Querbohrung des Störkörpers
\item DMS oder Piezo-Wandler an schwingungsfähigen Steg
\end{itemize}
Eigenschaften
\begin{itemize}
\item Fluideigenschaften: Gase, Dämpfe und/oder niedrigviskose Flüssigkeiten
\item Messgenaugikeit: $\pm 1 \%$ v.S.
\item Fehlerquellen: Schwingungen des Messrohres und Pulsation des Mediums
\end{itemize}
\begin{itemize}[label=$+$]
\item robust und wartungsarm
\item hohe Genauigkeit, großer Dynamikbereich
\item hohe Langzeitstabilität
\item Unempfindlichkeit gegen Verschmutzung
\item linearer Zusammenhang $f(v)$
\item beliebige Einbaulage
\end{itemize}
\begin{itemize}[label=$-$]
\item für kleine Durchflussgeschwindigkeiten nicht geeignet
\item relativ hohe verbleibende Druckverluste
\item große Ein- und Auslaufstrecke und ggf. Strömungsleichrichter nötig
\end{itemize}

\section{Magnetisch: Induktives Verfahren}
Abhängig von Lorentz-Kraft: $\originalVec{F}_L=q\cdot \originalVec{\nu} \times \originalVec{B}$\\
Spannung an dem Rohr: $U=B\cdot d \cdot \nu$\\
Eigenschaften:
\begin{itemize}
\item Medien: insbesondere wässrige Lösungen
\item Messgenaugikeit: $\pm 0,2\%$ bis $\pm 1\%$ v.S.
\item Fehlerquellen: Querschnittsverkleinerung durch Ablagerungen im Messrohr, relativ kleine Messsignale (im $\unit{mV}$-Bereich und niedriger)
\end{itemize}
\begin{itemize}[label=$+$]
\item berührungslos, robust, praktisch verschleiß- und wartungsfrei
\item unabhängig von Dichte, Temperatur, Druck, Viskosität des Mediums
\item unabhängig von Art der Strömung (laminar, turbulent) und vom Strömungsprofil
\item linearer Zusammenhang $q_v(U)$
\item keine Rohreinbauten notwending\\
$\to$  keine Querschnittsverengung $\to$ keine zusätzlichen Druckverluste
\item geeignet für korrosive, aggressive und verschmutzte Medien (Schlämme, Feststoffe)
\item kurze Ein- und Auslaufstrecken
\item Eignung für kleinste bis mittlere Mengen
\end{itemize}
\begin{itemize}[label=$-$]
\item nur für leitende Flüssigkeiten ($> 1 \unit{\mu S/cm}$) geeignet $\to$ nicht für Gase
\item es dürfen sich keine elektrisch leitenden Ablagerungen bilden
\item mit der Nennweitgröße überproportional steigende Herstellungskosten
\item die Elektroden sollten aus korrosionsfesten Material bestehen
\end{itemize}

\section{Ultraschall-Durchflussmessung}
Ultraschall durch Medium
\begin{itemize}
\item Diagonalverfahren: Messung diagonal zum Rohr mit zwei Sendern/Empfängern auf jeweils einer Seite des Rohrs.\\
 $\nu \approx \frac{c^2 \cdot \Delta t}{2 \cdot L \cdot \cos \varphi}$
\item Reflexionsverfahren: Messung mit einem Sender und einem Empfänger, durch Reflexion des Schalls am Rohr wird der Durchfluss gemessen.
\item Doppler-Verfahren: Beweglicher Sender/Empfänger. Durch Doppler-Effekt wird Durchfluss gemessen.\\
$\Delta f = 2 \cdot f_0 \cdot \cos \varphi \cdot \frac{\nu}{c}$
\end{itemize}
Eigenschaften:
\begin{itemize}
\item Medien: reine Flüssigkeiten und Gase (für Doppler: Streuzentren notwendig)
\item Messgenauigkeit: $\pm 0,5\%$ bis $3 \%$ v.S. (Doppler: schlechter)
\end{itemize}
\begin{itemize}[label=$+$]
\item keine Einbauten in der Rohrleitung\\
$\to$ keine Querschnittsverengung $\to$ keine zusätzlichen Druckverluste\\
Messungen über mehrere Meter Distanz möglich
\item linearer Messeffekt
\item Einsatz auch bei nichtleitenden und verschmutzten Fluiden (Doppler)
\item Clamp-On-Verfahren möglich
\end{itemize}
\begin{itemize}[label=$-$]
\item gut ausgebildete Strömungsprofile erforderlich
\item die Flüssigkeit muss frei von größeren Störkörpern sein (Feststoffe, Gasblasen)
\item die Laufzeitdifferenz ist Stoff- und temperaturabhängig
\item Schall schluckende bzw. reflektierende Ablagerungen sind zu vermeiden
\item relativ großer messtechnischer Aufwand
\item Beschränkung der Nennweitenpalette nach unten
\end{itemize}

\section{Korrelationsdurchflussmessung}
An Rohr werden (mit Abstand $d$) zwei Sensoren (für alles möglich: Druck, Temperatur usw.) befestigt. Diese werden beständig gemessen und erhalten dem entsprechend zwei Signale. Diese Signale werden auf Ähnlichkeiten untersucht. Der zeitliche Abstand dieser Ähnlichkeiten stellt den Durchfluss dar.\\
Eigenschaften:
\begin{itemize}
\item Fluideigenschaften: Flüssigkeiten, Gase, Mehrphasenströme
\item Messgenauigkeit: Abhängig von der Aufgabenstellung
\end{itemize}
\begin{itemize}[label=$+$]
\item dort, wo andere Verfahren versagen, besonders bei Mehrphasenströmungen
\item hohe Flexibilität bezüglich der Messstellen und Messstoffe
\item kann als Clamp-On-Messverfahren eingesetzt werden
\end{itemize}
\begin{itemize}[label=$-$]
\item manchmal sehr stark eingeschränkte Genauigkeit
\item aufgrund der rechnerischen Auswertung im Korrelator immer langsamer als direkt anzeigende Messverfahren
\end{itemize}

\section{Hitzdraht-Anemometer}
Kalorimetrie (calometry): Messung von Wärmemengen
\subsection{Hitzdraht-Anemometrie}
Messung der geschwindigkeitsabhängigen Wärmeabgabe eines beheizten Sensors an ein vorbeifließendes, kälteres Fluid.
$$q_m = A \left( \frac{R_S \cdot I^2}{a\cdot (T_S-T_0)}-\frac{b}{a}\right)$$
wobei $a,b$ und $A$ Konstanten sind.\\
konstanter Heizstrom CC (langsam): $I=const \; R=R(\nu)$\\
konstante Hitzdrahttemperatur CT (schneller): $R=const \; I=I(\nu)$
\begin{itemize}
\item Fluideigenschaften: hauptsächliche Gase
\item Messgenauigkeit: $(\pm 2) \pm 1 \%$ v.S.
\end{itemize}
\begin{itemize}[label=$+$]
\item empfindlich, besonders geeignet für kleine und mittlere Strömungsgeschwindigkeiten
\item geringe Abmessung des Heizdrahts $\rightarrow$ kleine Einstellungszeit
\item rel. geringe Heizleistung erforderlich
\item wenige Druckverluste an der Messstelle
\item direkte elektrische Anzeige des Messwertes
\item sowohl im lam. und turb. Strömungen (Kalibrierung notwendig)
\end{itemize}
\begin{itemize}[label=$-$]
\item abnehmende Empfindlichkeit mit wachsender Geschwindigkeit
\item Messung bei kleinen Durchflüssen wegen Eigenkonvektion nicht möglich
\item Eichung jedes einzelnen Drahtes notwendig
\item Alterung infolge von temperaturbedingten Strukturänderungen
\item abhängig vom Medium und Strömungsart
\end{itemize}
\subsection{Aufheiz-Anemometrie}
Messung der Erwärmung des zu messenden Fluids durch Wärmestromzufuhr aus einer Heizwicklung.
$$q_m=f(\Delta T)$$
\begin{itemize}
\item hohe Heizleistung erforderlich
\item linearer Messeffekt
\item unabhängig von Dichte des Mediums
\item abhängig von spezifischer Wärme
\item Genaugikeit: $\leq 1 \%$ v.E.
\item Einstellzeit von Null bis $\pm 2 \%$ v.S.: $< 2\unit{s}$
\end{itemize}

\section{Coriolis Durchflussmessung}
\begin{itemize}
\item inertiales Bezugssystem
\item rotierendes Bezugssystem
\end{itemize}
$\vec{F}_C=2m\cdot (\vec{v}\times \vec{\omega})$\\
Rechte Hand: Daumen $\vec{F}_C$, Zeigefinger $\vec{v}$, Mittelfinger $\vec{\omega}$ (gegen Uhrzeigersinn drehend)\\
$q_m \sim \varphi$

\begin{itemize}
\item Fluideigenschaften: Flüssigkeiten und Gase bei höheren Drücken ($>5\unit{bar}$)
\item Breie und ähnliche hochviskose Fluide, homogene Mehrphasengemische
\item pneumatisch förderbare Schüttgüter
\item Messgenauigkeit: $\pm 0,2 - 0,3\%$ v.S. für Flüssigkeiten und Gase, $\pm 1\%$ v.S. für pneumatisch förderbare Schüttgüter
\end{itemize}
\begin{itemize}[label=$+$]
\item Eine direkte Bestimmung des Massenstroms
\item nahezu unabhängig von Druck, Temperatur, Dicht, Viskosität des Mediums
\item Resonanzfrequenz $\sim$ Dichte des Mediums
\item auch bei Kleinstmengen, pulsierende Strömung, nicht vollständig befüllten Rohren, hohen und niedrige Temperaturen und hohe Drücken einsetzbar
\end{itemize}
\begin{itemize}[label=$-$]
\item hohe Druckverluste
\item eine Schwingungsentkopplung der Messanlage erforderlich
\item empfindlich zu Druckschwankungen
\item verklebende Schüttgüter können das Messergebnis verfälschen
\item unsicher bei niedrigen Strömungsgeschwindigkeiten von Gasen
\end{itemize}

\chapter{Binäre Positionssensoren}
nur zwei Ausgangssignale (Schaltzustände): „Ein“/„Aus“
\begin{itemize}
\item Taktile Sensoren (Taster):\\
mit mechanischem Kontakt mit dem zu detektierenden Objekt
\item Näherungsschalter (Initiatoren):\\
berührungslos arbeitende Schalter
\end{itemize}
\begin{tabular}{l | l}
Sensortyp & berührungslos\\\hline
Grenztaster & nein \\
Lichtschranken & ja \\
Induktive Sensoren & ja\\
Kapazitive Sensoren & ja\\
Ultraschall-Sensoren & ja\\
(Passiv-)Infrarotmelder & ja
\end{tabular}

\paragraph{Näherungsschalter (gegen Taster)}
\begin{itemize}[label=$+$]
\item berührungslos
\item rückwirkungsfrei
\item verschleißfrei
\item erzeugen prellfreie Ausgangssignale
\item z.T. über große Entfernungen
\item Lebensdauer $>$ als bei Tastern
\end{itemize}
\begin{itemize}[label=$-$]
\item Hilfsenergie notwendig
\item Fehlschaltungen durch Fremdfelder und Verschmutzung
\item z.T. wesentlich teurer als Taster
\end{itemize}

\section{Grenztaster}
Messprinzip: Kontakt über Stößel betätigen (taktil) [Schaltkraft]\\
Ausführungen: gekapselt/leicht\\
Schaltpunktgenauigkeit $\pm 10 \unit{\mu m}$
\subsection{Präzisionsschalter}
\begin{tabular}{l l}
Schaltpunktgenauigkeit & $\pm 1 \unit{\mu m}$\\
Hysterese & $< 2 \unit{\mu m}$\\
mech. Lebensdauer & $10^6$\\
Schaltfrequenz & $10 \unit{Hz}$\\
Nachlaufweg Stössel & $1,5\unit{mm}$
\end{tabular}
\subsection{Kontaktprellen}
Prellen: das Schaltelement bei einem Schließvorgang durch die Federwirkung mehrmals öffnet oder wieder schließt, bis ein endgültiger Kontakt da ist.
\begin{itemize}
\item idealer Schalter: kein Zeitverzug (vertikale Flanke)
\item Näherungsschalter: schaltet im $\unit{\mu s}$-Bereich (Flanke mit Anstieg)
\item Mechanischer Schalter: höhere Verzögerung schaltet im $\unit{ms}$-Bereich zusätzlich Kontaktprellen (Schalter federt nach, es gibt mehrere kleinere Flanken in einem Bereich)
\end{itemize}
\subsection{Vor- und Nachteile von Grenztastern}
\begin{itemize}[label=$+$]
\item Preisgünstig: viel preisgünstiger als Initiatoren
\item Robust und sicher im Einsatz
\item klein in den Abmessungen
\item keine Hilfsenergie notwendig
\item Verwendbarkeit bis $600 \unit{V}$
\item unbeeinflussbar durch Fremdfelder
\item sehr hohe Genauigkeit und Wiederholgenauigkeit des Schaltpunktes (typisch $\pm 10 \unit{\mu m}$) auch bei low-cost Typen
\item für kleinste und größte Schaltleistungen erhältlich
\item Absolut zuverlässige (hundertprozentige) galvanische Trennung
\end{itemize}
\begin{itemize}[label=$-$]
\item eine mechanische Berührung erforderlich (Schaltkraft, Eingriff in Prozess)
\item spez. mecha. Vorrichtungen notwendig
\item Kontakte prellen $\to$ unsauberes Signal
\item hohe Anfälligkeit gegen Verschmutzung und Verschleiß
\item nur niedrige Schaltfrequenzen realisierbar
\item Schaltpunktdrift durch Verschleiß und Materialermüdung
\item Übergangswiderstände an den Kontakten
\end{itemize}
Messgrößen:
\begin{itemize}
\item Objektdistanz über/unterschritten
\item Druck, Kraft
\item Niveau
\end{itemize}
Einsatzgebiete:
\begin{itemize}
\item End- oder Referenzschalter
\item Sicherheitsschalter
\item Ablaufsteuerungen
\item absolute Nullpunktgeber eines Inkrementalgebers
\item Präzisionsmessung
\item …
\end{itemize}
%\newpage
%\printbibliography

\end{document}
